% Source presentations for register row N-c-017.

\subsection{Four mechanisms behind symbolic verification}
\label{sec:p09-presentation-n-c-017}

\subsubsection{Bryant's fixed-order reduced function graphs}
\label{sec:p09-n-c-017-bryant}

This treatment follows Bryant's annotated electronic author version;
section and figure locators refer to that version \cite{Bryant1986}.

Fix \(n\in\mathbb N\) and Boolean arguments in the order
\(x_1<\cdots<x_n\). A Boolean function has type
\(f:\{0,1\}^n\to\{0,1\}\). For \(1\leq i\leq n\) and
\(b\in\{0,1\}\), its cofactor is the function
\[
 (f|_{x_i=b})(a_1,\ldots,a_n)
 :=f(a_1,\ldots,a_{i-1},b,a_{i+1},\ldots,a_n).
\]
It is independent of its \(i\)-th argument. At \(n=0\), the domain contains
the single empty assignment and the two functions are constants.
Bryant's
\emph{function graph} is a rooted directed graph in which every vertex is
reachable from the root.  A terminal vertex has a value in \(\{0,1\}\).  A
nonterminal vertex \(v\) has an index \(i\) and distinguished children
\(\operatorname{low}(v)\) and \(\operatorname{high}(v)\); every nonterminal
child has index strictly greater than \(i\).  The ordering condition makes the
graph acyclic.  If \(f_v\) is the Boolean function denoted by the subgraph
rooted at \(v\), then
\[
 f_v=
 \begin{cases}
  \operatorname{value}(v),&v\text{ terminal},\\
  (\neg x_i\land f_{\operatorname{low}(v)})
   \lor(x_i\land f_{\operatorname{high}(v)}),
   &\operatorname{index}(v)=i.
 \end{cases}
\]
Thus an assignment follows the low edge when \(x_i=0\), the high edge when
\(x_i=1\), and obtains the terminal value at the end of its unique path
\cite[defs.~1--2]{Bryant1986}.

The graph is \emph{reduced} exactly when it has neither a vertex with identical
low and high children nor two distinct vertices whose rooted subgraphs are
isomorphic.  The first forbidden pattern is an irrelevant test; the second is
duplicate representation of one residual function.  Eliminating the former
and merging the latter preserve the denoted function
\cite[defs.~3--5]{Bryant1986}.

For example, expand \(x_1\land x_2\) as the full ordered decision tree.
First merge its three \(0\)-leaves; there is only one \(1\)-leaf.
Below \(x_1=0\), both edges of the \(x_2\)-vertex now reach that same
\(0\)-terminal, so the identical-child reduction removes the test.
The reduced
graph has a root testing \(x_1\), low child \(0\), and high child testing
\(x_2\) with children \(0\) and \(1\).  Both reductions are needed: one removes
an irrelevant decision, while the other shares equal residual functions.

For a Boolean function \(f\), put
\[
 I_f=\{i:f|_{x_i=0}\ne f|_{x_i=1}\}.
\]
Bryant's Theorem~1 says that \(f\) has a unique reduced function graph up to
isomorphism and that every nonisomorphic function graph denoting \(f\) under
the same order has more
vertices \cite[thm.~1]{Bryant1986}.

For existence, start with the full ordered decision tree of \(f\).
Each reduction preserves the function and order and strictly decreases
the number of vertices. Retain only vertices reachable from the new
root after each step. This finite process ends in a reduced graph.

Prove uniqueness by induction on the number of variables still
allowed in an ordered suffix \(x_j,\ldots,x_n\). The empty suffix
allows only a terminal, fixed by the constant function's value.
For the induction step, the two children of any root testing
\(x_\ell\), \(\ell\geq j\), use only the strictly shorter suffix
\(x_{\ell+1},\ldots,x_n\). Hence the induction hypothesis already
applies to their reduced subgraphs, whether or not \(x_\ell\)
is essential to the function.

If \(f\) is constant, the two children of a nonterminal root would
denote that same constant. By induction they would be isomorphic,
and reducedness would make them the same vertex. The root would
then be a forbidden identical-child test. Thus a reduced constant
graph is its single terminal.
Otherwise put \(i=\min I_f\). A root index greater than \(i\)
would skip an essential variable. A root index \(\ell<i\)
would have equal cofactor functions because \(f\) is independent
of \(x_\ell\). The shorter-suffix induction would again identify
the child subgraphs, contradicting reducedness. Therefore the
root index is exactly \(i\), and its low and high functions are
the two prescribed cofactors at \(x_i\).

Induction determines both reduced child graphs up to isomorphism.
These isomorphisms also determine their sharing: whenever vertices
in the low and high graphs denote the same residual function,
shorter-suffix uniqueness makes their rooted graphs isomorphic,
and reducedness requires one common vertex. Conversely, a shared
vertex necessarily has one denotation. Thus the two child
isomorphisms agree on their overlaps and extend to an isomorphism
of the whole graphs. This completes the induction on suffix length.

Finally, a nonisomorphic graph for \(f\) cannot already be reduced.
Reducing it therefore makes at least one strict size decrease
before reaching the unique reduced graph. This proves minimality.
The theorem is conditional on the one fixed argument order;
it compares no two orders
\cite[thm.~1]{Bryant1986}.

\paragraph{Bottom-up reduction.}
Let \(N=|G|\) count the vertices of the input graph, including its terminals.
Assign terminals the bookkeeping index \(n+1\). The reduction procedure
assigns each original vertex \(v\) an integer \(\operatorname{id}(v)\) for
its reduced residual function, and selects one representative vertex for
each identifier. Figure~4 carries out the following operations
\cite[\S4.2 and fig.~4]{Bryant1986}.

\begin{enumerate}
\item Collect the reachable vertices by index and process the levels from
\(n+1\) down to \(1\). Every child of a vertex at the current nonterminal
level has already been processed.
\item At the terminal level, use the value \(0\) or \(1\) as key. All
terminals with the same value receive one identifier and one representative.
\item At a nonterminal vertex \(v\), compare the identifiers of its two
children. If they are equal, set
\(\operatorname{id}(v)=\operatorname{id}(\operatorname{low}(v))\);
this vertex contributes no new representative.
\item For the remaining vertices at level \(i\), sort the pairs
\[
 \bigl(\operatorname{id}(\operatorname{low}(v)),
       \operatorname{id}(\operatorname{high}(v))\bigr).
\]
Give equal pairs one identifier, select one representative for each distinct
pair, and redirect its children to their selected representatives. Different
levels are processed separately, so equal pairs at different levels are not
merged.
\item Return the representative with the original root's identifier, retaining
the subgraph reachable from it.
\end{enumerate}

The correctness invariant can be proved from the two reductions. After a
level is processed, every processed vertex denotes the same function as its
representative, and the representative subgraph is reduced. At a terminal
the invariant follows from its value. Suppose it holds below level \(i\).
If the child identifiers agree, the two cofactors agree, and Shannon
expansion reduces the current function to that common cofactor. Otherwise,
the ordered child pair determines the function at \(x_i\); merging equal
pairs preserves it. Because the two reduced children have different
functions, the result depends on \(x_i\), so it cannot equal a representative
whose first essential variable has a different index. Thus no redundant
test or duplicate residual remains. Induction reaches the root and gives the
unique reduced graph from Theorem~1.

The source's Figure~5 makes both operations visible. Its four terminal
records have values \(0,1,0,1\), receiving identifiers \(1,2,1,2\).
At index \(3\), both child-pair keys are \((1,2)\), receiving identifier \(3\).
At index \(2\), key \((1,3)\) receives \(4\), while key \((3,3)\) is eliminated
and inherits \(3\). The root's key is then \((4,3)\), receiving \(5\).
The output has one representative for each of the five identifiers.

Comparison sorting gives
\[
 \sum_i O(m_i\log(1+m_i))=O(N\log(1+N)),
 \qquad \sum_i m_i\leq N,
\]
where \(m_i\) counts the keyed vertices at level \(i\). If an implementation
scans an array of all \(n+1\) levels, its bookkeeping also costs \(O(n)\).
Thus this comparison-sorting implementation runs in
\[
 O\bigl(n+N\log(1+N)\bigr).
\]
With a reduction implementation that processes each vertex in constant
amortized or expected time, the graph-processing cost is \(O(N)\), under
that implementation's data-structure assumptions. Either way, reduction
of the intermediate result is part of the cost of \textsc{Apply} and
\textsc{Compose} \cite[\S4.2 and table~1]{Bryant1986}.

\paragraph{Residual functions and the effect of variable order.}
Fix a cut in the ordered variable list. For any assignment to the
variables before the cut, follow its graph edges until a terminal
or the first test after the cut. The reached subgraph denotes the
residual function of the remaining variables. Distinct residual
functions therefore require distinct reached vertices, even if
some input tests are skipped. Thus \(L\) distinct prefix residuals
force at least \(L\) vertices in any ordered graph, reduced or not.

The source's ordering example makes this principle quantitative. Under
the displayed argument order,
\[
 \begin{aligned}
  g_3&=x_1x_2\lor x_3x_4\lor x_5x_6,\\
  h_3&=x_1x_4\lor x_2x_5\lor x_3x_6
 \end{aligned}
\]
have reduced graphs with respectively eight and sixteen vertices, although
one formula is obtained from the other by permuting arguments.
For \(n\geq1\), on \(2n\) Boolean arguments,
\(\bigvee_{r=1}^{n}x_{2r-1}x_{2r}\) has \(2n+2\) vertices, whereas
\(\bigvee_{r=1}^{n}x_rx_{n+r}\) requires \(2^{n+1}\)
\cite[\S3.2 and fig.~2]{Bryant1986}.

Here is the exact count. For the interleaved expression, put
\[
 H_j=\bigvee_{r=j}^{n}(x_{2r-1}\land x_{2r}),\qquad H_{n+1}=0.
\]
Use an \(x_{2j-1}\)-vertex with low child \(H_{j+1}\) and high
child an \(x_{2j}\)-vertex; that second vertex has low child
\(H_{j+1}\) and high child \(1\). All \(2n\) test vertices and
both terminals are reachable. Neither type of test has equal
children: make every later pair false and vary the current
pair to distinguish the branches. The test indices are all
distinct, so there are no duplicate rooted subgraphs. This
is a reduced graph with exactly \(2n+2\) vertices.

For the separated expression, write \(y_j=x_{n+j}\).
Before testing \(x_i\), the already chosen prefix selects a subset
\(A\subseteq\{1,\ldots,i-1\}\), leaving
\[
 R_{i,A}=\bigvee_{j\in A}y_j
       \lor\bigvee_{j=i}^{n}(x_j\land y_j).
\]
Different \(A\)'s give different functions, as an assignment
making just one \(y_j\) true distinguishes any differing
membership. Every \(R_{i,A}\) depends on \(x_i\): set \(y_i=1\)
and every other \(y\)-variable to zero. Hence there are
\(2^{i-1}\) distinct roots at that level, and the first block
contains \(\sum_{i=1}^{n}2^{i-1}=2^n-1\) vertices.
After this block, every subset \(A\subseteq\{1,\ldots,n\}\)
occurs and leaves \(Q_A=\bigvee_{j\in A}y_j\).
The empty subset gives \(0\). Each nonempty subset gives a
distinct root testing \(y_{\min A}\), with high child \(1\)
and low child \(Q_{A\setminus\{\min A\}}\). These are exactly
the \(2^n-1\) nonterminal residuals in the second block.
Adding the two terminals gives \(2^{n+1}\).
At \(n=0\), both empty disjunctions are constant \(0\) and
their graphs have one vertex; the positive-\(n\) formulas
do not apply.

\paragraph{Symmetry can keep the graph small.}
Odd parity on \(n\geq1\) variables has one residual at its
first level and two at every later level: parity of the
unread variables or its complement, according to the
prefix parity. Both later residuals occur, are distinct,
and depend on the next variable. With the two terminals,
the reduced graph therefore has
\[
 1+2(n-1)+2=2n+1
\]
vertices, in any input order. In contrast, a noncontradictory
conjunction of literals
implying odd parity cannot omit an input variable, since
flipping an omitted variable changes parity. Every
noncontradictory such conjunction therefore specifies a
full satisfying assignment. There are \(2^{n-1}\) of those,
so that many full minterms are necessary and sufficient
in a disjunctive normal form.

More generally, if
\(f(x_1,\ldots,x_n)=F(x_1+\cdots+x_n)\) for
\(F:\{0,\ldots,n\}\to\{0,1\}\), a prefix of \(k\) tests
needs only its number \(t\) of ones. Store a state \((k,t)\);
the next branches lead to \((k+1,t)\) and \((k+1,t+1)\),
and after the final test output \(F(t)\).
For \(n\geq1\), merging equal terminal values and reducing
therefore gives
\[
 |G_f|\leq\sum_{k=0}^{n-1}(k+1)+2
          =\frac{n(n+1)}2+2.
\]
This is an upper bound, not an asserted exact size for
every symmetric function. The construction works in
every variable order because the number of ones is
permutation-invariant \cite[\S3.1]{Bryant1986}.

\paragraph{Combining two graphs.}
The source's \textsc{Apply} procedure computes \(f_1\mathbin{\circ}f_2\) for
any binary Boolean operation \(\circ\).  At roots \(v_1,v_2\), it chooses the
smaller current index \(i\), recursively combines the two low restrictions and
the two high restrictions, and records the result for the pair
\((v_1,v_2)\).  If only one root tests \(x_i\), the other root is reused in both
recursive calls.  A controlling terminal value, such as \(1\) for disjunction,
ends that branch immediately.  Memoization ensures that each vertex pair is
expanded once; reduction then removes redundant tests and duplicate subgraphs.

Figure~7 applies disjunction to
\[
 f_1=\neg(x_1x_3),
 \qquad
 f_2=x_2x_3.
\]
The pair consisting of the first graph's terminal \(1\) and the second graph's
root stops by the controlling-value rule; a later repeated pair is returned
from the memo table.  Reducing the evaluation graph yields
\[
 f_1\lor f_2=\neg(x_1\land\neg x_2\land x_3).
\]
Put \(P=|G_1||G_2|\). There are at most \(P\) vertex pairs and at most \(P\)
vertices in the unreduced evaluation graph. With constant-time memo-table
access, recursive work is \(O(P)\); comparison-sorting reduction gives total
time \(O(n+P\log(1+P))\). A linear reduction implementation gives \(O(P)\)
graph-processing time under its corresponding access assumptions.

The product dependence can be necessary for the output itself.
Take \(n,m\geq1\) and use the order
\[
 x_1<\cdots<x_{n+m}<y_1<\cdots<y_{n+m}.
\]
Let
\[
 f_1=\bigvee_{i=1}^{n}(x_i\land y_i),\qquad
 f_2=\bigvee_{i=n+1}^{n+m}(x_i\land y_i).
\]
Ignoring variables on which an operand does not depend, each
has the separated-pair order already counted. Thus
\[
 |G_1|=2^{n+1},\qquad |G_2|=2^{m+1}.
\]
Their disjunction is the separated-pair function on all
\(n+m\) pairs, so its reduced graph has
\[
 |G_{f_1\lor f_2}|=2^{n+m+1}
                 =\tfrac12|G_1||G_2|
\]
vertices. Any method materializing that ordered output graph
must therefore perform work proportional to the product
on this family \cite[\S4.3]{Bryant1986}.

This remains an input-graph bound. The multiplication argument below
proves that, for every variable order and \(n\geq1\), some output bit of
\(n\)-bit integer multiplication needs at least \(2^{n/8}\) vertices
\cite[\S\S4.3, 6 and app.]{Bryant1986}.

\paragraph{Restriction, quantification, and equality.}
Restriction redirects an \(x_i\)-test to its selected child and reduces the
result \cite[\S4.4]{Bryant1986}. The following identities make the subsequent
symbolic operations explicit:
\[
\begin{aligned}
 \exists x_i\,f&=f|_{x_i=0}\lor f|_{x_i=1},\\
 \forall x_i\,f&=f|_{x_i=0}\land f|_{x_i=1},\\
 f|_{x_i=g}
   &=(\neg g\land f|_{x_i=0})\lor(g\land f|_{x_i=1}).
\end{aligned}
\]
In the last equation \(g:\{0,1\}^n\to\{0,1\}\) is evaluated on the original
assignment, including its original \(x_i\)-value if it depends on \(x_i\).
The equation is therefore simultaneous functional substitution, not a
recursive equation for \(x_i\). Each identity follows by the two possible
truth values of the quantified or substituted argument.

For \(W\subseteq\{x_1,\ldots,x_n\}\), eliminate its members one at a
time. For existential abstraction, any order gives the same function because
\[
 (\exists W\,f)(a)
 =\bigvee_{b\in\{0,1\}^{W}} f(a[W:=b]).
\]
Here \(a[W:=b]\) agrees with \(b\) on \(W\) and with \(a\) elsewhere.
If \(W=\varnothing\), there is one empty assignment \(b\), so abstraction
returns \(f\). Abstraction order may still affect intermediate graph sizes.

For example, \(f(x,y)=x\land y\) has cofactors
\(f|_{y=0}=0\) and \(f|_{y=1}=x\). Hence
\(\exists y\,f=x\), whereas \(\forall y\,f=0\).
This calculation derives a set projection directly from two restrictions.

Canonicity makes equality of reduced graphs under one order an equality
decision for their functions. Constant-time root comparison additionally
requires a common representation in which equal reduced subgraphs share one
identifier; separately allocated isomorphic graphs need not have the same
root address. Consequently the symbolic fixed-point stopping test below uses
semantic BDD equality, or identifier equality under that sharing convention.

\paragraph{Composition by three residuals.}
For Boolean functions \(a,b,c\), define
\[
 \operatorname{ITE}(a,b,c):=(a\land b)\lor(\neg a\land c).
\]
Composition \(f_1|_{x_i=f_2}\) is
\[
 \operatorname{ITE}(f_2,f_1|_{x_i=1},f_1|_{x_i=0}).
\]
Figure~8 combines restriction and this ternary operation in one traversal.
Its recursive state consists of a low-cofactor vertex from \(G_1\), a
high-cofactor vertex from \(G_1\), and a vertex from \(G_2\). Whenever a
cofactor vertex tests \(x_i\), replace it by the corresponding low or high
child before looking up the triple. At a terminal triple, evaluate
\(\operatorname{ITE}\). Otherwise split at the smallest current index,
retain an unchanged vertex in both branches when it does not test that index,
recurse on the resulting low and high triples, and memoize the result.
Reduce the returned graph \cite[\S4.5 and fig.~8]{Bryant1986}.

For a fixed assignment, the cofactor replacements fix only the selected
argument of \(f_1\). The third residual still evaluates \(f_2\) on the original
assignment. Each split partitions assignments by one variable, so induction
on the remaining indices proves that the returned function is the displayed
\(\operatorname{ITE}\). There are at most
\(|G_1|^2|G_2|\) triples; this bounds recursive pairings before the separately
charged reduction cost. For \(C=|G_1|^2|G_2|\), comparison-sorting reduction
therefore gives \(O(n+C\log(1+C))\) total time; a linear reduction routine
gives \(O(C)\) graph-processing time under its access assumptions.

For example, let \(f_1=x_1\mathbin{\oplus}x_3\) and
\(f_2=x_2\), and substitute for \(x_1\). The two cofactors are
\(f_1|_{x_1=0}=x_3\) and \(f_1|_{x_1=1}=\neg x_3\); hence
\[
 f_1|_{x_1=x_2}
   =\operatorname{ITE}(x_2,\neg x_3,x_3)
   =x_2\mathbin{\oplus}x_3.
\]
The example also shows why a functional substitution may change the first
tested variable.

Graph splicing is another implementation of substitution when the resulting
paths respect the fixed variable order. After splicing, eliminate
identical-child tests and merge isomorphic residual subgraphs. Preserving
variable order alone does not ensure either reduction property.

\paragraph{Finding, listing, and counting satisfying assignments.}
Let
\[
 S_f:=\{a\in\{0,1\}^n:f(a)=1\}.
\]
For a reduced graph, Lemma~3 states that every nonterminal vertex has a
\(1\)-terminal descendant \cite[\S4.6]{Bryant1986}. Indeed, if all its
terminal descendants were \(0\), its function would be constant \(0\), whose
unique reduced graph has no nonterminal vertex.

To find one assignment, initialize the \(n\) argument bits arbitrarily.
At a terminal return failure for \(0\) and success for \(1\). At an
\(x_i\)-test, first set \(x_i=0\) and try the low child; on failure set
\(x_i=1\) and try the high child. Reducedness and Lemma~3 imply that a failing
first choice goes immediately to terminal \(0\); the other choice succeeds.
Thus the search follows at most \(n\) tests and outputs a full assignment in
\(O(n+1)\) time, including the constant-function cases.

To list every assignment, keep the next output index \(i\) as well as the
current vertex \(v\). When \(\operatorname{index}(v)>i\), branch over both
values of the skipped variable while retaining \(v\). When the indices
agree, follow the low and high children for the corresponding bit choices.
Discard a \(0\)-terminal immediately; at \(i=n+1\), output the completed
assignment exactly when the terminal is \(1\). The invariant is that the
already chosen prefix is fixed and \(f_v\) describes precisely its possible
satisfying extensions. Every assignment is generated once. For \(n\geq1\)
and \(S_f\ne\varnothing\), the output-sensitive cost is \(O(n|S_f|)\);
an empty satisfying set needs only the constant-time root test.

Counting avoids enumeration. For a vertex \(v\) of index \(i\), let
\(\alpha_v\) count satisfying assignments to variables \(x_i,\ldots,x_n\)
for its residual function. A terminal has index \(n+1\) and
\(\alpha_v=\operatorname{value}(v)\). The recurrence is
\[
 \alpha_v=
  2^{\operatorname{index}(\operatorname{low}(v))-i-1}
       \alpha_{\operatorname{low}(v)}
  +2^{\operatorname{index}(\operatorname{high}(v))-i-1}
       \alpha_{\operatorname{high}(v)}.
\]
For root \(r\),
\[
 |S_f|=2^{\operatorname{index}(r)-1}\alpha_r.
\]
Each child edge skips precisely
\(\operatorname{index}(\text{child})-i-1\) untested variables, each freely
Boolean; the branch fixes \(x_i\) itself. The two branches are disjoint.
These facts prove the recurrence, and the root factor restores the variables
skipped before the first test.

For \(f=x_1\) on one argument, the accepting edge skips no variables and
contributes \(2^0=1\); the rejecting edge contributes \(0\).
For the constant \(1\), the root factor gives \(2^n\), including one empty
assignment when \(n=0\).

A further check uses \(f=x_1\lor x_3\) on three arguments. The \(x_3\)-vertex
has \(\alpha=1\). The root's low edge skips \(x_2\), contributing \(2\);
its high edge reaches terminal \(1\), leaving \(x_2,x_3\) free and
contributing \(4\). Hence \(|S_f|=6\). A bottom-up traversal evaluates the
recurrence once per vertex. The paper's \(O(|G|)\) bound assumes
constant-time integer arithmetic at sufficient precision; the counts may
require \(n+1\) bits, so this is not a bit-complexity bound.

\paragraph{An exponential lower bound for multiplication.}
Let \(n\geq1\) and define unsigned inputs
\[
 a=\sum_{u=1}^{n}a_u2^{u-1},\qquad
 b=\sum_{v=1}^{n}b_v2^{v-1}.
\]
Write \(\operatorname{mul}_r(a,b)\) for bit \(r\) of the full
\(2n\)-bit product, numbering bits from \(1\) at the least
significant end. For every order of the \(2n\) input variables,
there exists \(r\in\{1,\ldots,2n-1\}\) such that every ordered
graph for that output has at least
\[
 2^{q_*}\text{ vertices},\qquad
 q_*=\left\lceil
       \frac{\lceil n/2\rceil^2}{2n-1}
      \right\rceil.
\]
In particular its size is at least \(2^{n/8}\). The selected
output bit may depend on the input order
\cite[app., thm.~2 and lem.~4]{Bryant1986}.

\emph{Proof.}
Cut the reading order after its first \(n\) variables.
If it has \(t\) bits of \(a\) before the cut, it has \(t\)
bits of \(b\) after it. Choose \(a\) as a data operand and
\(b\) as a control operand when \(t\geq n/2\); otherwise
exchange their roles. Multiplication is commutative.
Let \(A\) be the significance indices of data bits before
the cut and \(B\) those of control bits after it. Then
\[
 |A|=|B|=p=\max\{t,n-t\}\geq\lceil n/2\rceil.
\]
These are bit-significance indices, not positions in the
reading order.

For \(1\leq r\leq2n-1\), count the pairs on an output diagonal:
\[
 q_r=|\{(u,v)\in A\times B:u+v-1=r\}|.
\]
Every pair belongs to exactly one diagonal, so
\(\sum_rq_r=p^2\). Some \(r\) therefore satisfies
\[
 q_r\geq\left\lceil\frac{p^2}{2n-1}\right\rceil
       \geq q_*.
\]
Let \(A_r=\{u\in A:r-u+1\in B\}\); the corresponding
control index is unique, so \(|A_r|=q_r\).
For each vector \(\alpha\in\{0,1\}^{A_r}\), assign those
data bits the values \(\alpha_u\), and assign every other
variable before the cut zero.

Two distinct such prefixes differ at some \(u\in A_r\).
Give both the same suffix: set the control bit at
\(v=r-u+1\in B\) to \(1\), and every other variable after
the cut to zero. The complete control operand is
\(2^{v-1}\), so multiplication is an exact left shift and
\[
 \operatorname{mul}_r=\operatorname{data}_{r-v+1}
                    =\operatorname{data}_u.
\]
The two outputs differ. Thus the \(2^{q_r}\) prefixes
induce pairwise different residual functions at the cut.
The residual principle forces at least \(2^{q_r}\)
vertices, even without assuming reducedness.
Finally,
\[
 \frac{p^2}{2n-1}
  \geq\frac{n^2}{4(2n-1)}>\frac n8,
\]
which yields the stated weaker exponential bound as well.
For \(n=1\), the construction has one selected data bit,
one control bit, and \(r=1\), so no small-\(n\) exception
is needed. \(\square\)

The quantifiers do not assert that every output is hard:
the least significant output is just \(a_1\land b_1\).
They assert that no choice of one input order makes all
product outputs small.

\paragraph{The ALU experiment.}
Bryant also compared gate-level and behavioural descriptions
of arithmetic logic units assembled from 74181 and 74182
components. Table~2 reports the following VAX~11/780 CPU
times for constructing both descriptions' graphs and checking
their equivalence \cite[\S5 and table~2]{Bryant1986}.
\begin{center}
\begin{tabular}{rrrr}
\hline
Word bits & Gates & CPU minutes & Largest output graph\\
\hline
4  & 52  & 1.1  & 197\\
8  & 123 & 2.3  & 377\\
16 & 227 & 6.3  & 737\\
32 & 473 & 22.8 & 1457\\
64 & 927 & 95.8 & 2897\\
\hline
\end{tabular}
\end{center}
The largest graph in each row is the equality-indication output.
With two \(n\)-bit operands, five controls, and carry-in,
the input space has \(2^{2n+6}\) patterns; the experiment
does not enumerate them. The best tested order puts the
controls first, then carry-in, then interleaves operand
bits from least to most significant. It is a measured
choice, not a proof of global ordering optimality.
Figure~11 contrasts these orders with separated operands,
illustrating how the chosen representation affects graph
growth in this particular circuit family
\cite[fig.~11]{Bryant1986}.

\subsubsection{Burch et al.'s symbolic fixed points}
\label{sec:p09-n-c-017-burch}

The paper's primary specification language is a relational \(\mu\)-calculus;
its symbolic algorithm evaluates it over finite domains. Its formulas are
built from applications
\(R(z_1,\ldots,z_r)\), Boolean connectives, and existential quantification over
individual variables.  An \(r\)-ary relational term is a relational variable,
an abstraction \(\lambda z_1,\ldots,z_r[f]\), or a least fixed point
\(\mu P[R]\).  The last form is admitted only when every free occurrence of
\(P\) in \(R\) lies under an even number of negations, so that the induced
operator is monotone.  The greatest fixed point \(\nu P[R]\) is introduced by
duality. These terms denote relations over one nonempty domain \(D\); the
fixed-point order is inclusion
\cite[\S3]{BurchEtAl1992}.

\paragraph{The relation denoted by a term.}
Take \(D\ne\varnothing\). Each relational variable has a positive arity \(r\);
an interpretation \(I_R\) assigns it an element of \(\mathcal P(D^r)\),
and \(I_D\) assigns individual variables elements of \(D\).
An abstraction binds pairwise distinct variables. In an application
\(R(z_1,\ldots,z_r)\), the actual argument variables are not free in \(R\);
this is a source syntax condition, not a claim that \(R\) has no other free
individual parameters \cite[\S3]{BurchEtAl1992}.

The following denotation clauses spell out the two sorts. Write
\(\llbracket R\rrbracket_{I_R,I_D}\subseteq D^r\) for a relational term,
and \(I_R,I_D\models f\) for a formula:
\[
\begin{aligned}
 \llbracket P\rrbracket_{I_R,I_D}&=I_R(P),\\
 \llbracket\lambda z_1,\ldots,z_r[f]\rrbracket_{I_R,I_D}
 &=\{(a_1,\ldots,a_r)\in D^r:
       I_R,I_D[z_1:=a_1,\ldots,z_r:=a_r]\models f\},\\
 I_R,I_D\models R(z_1,\ldots,z_r)
 &\iff (I_D(z_1),\ldots,I_D(z_r))
             \in\llbracket R\rrbracket_{I_R,I_D},\\
 I_R,I_D\models\neg f&\iff I_R,I_D\not\models f,\\
 I_R,I_D\models f\lor g
 &\iff (I_R,I_D\models f)\text{ or }(I_R,I_D\models g),\\
 I_R,I_D\models\exists z[f]
 &\iff\exists a\in D\;(I_R,I_D[z:=a]\models f).
\end{aligned}
\]
Conjunction and universal quantification are the Boolean duals of disjunction
and existential quantification. Interpretation updates change only the named
variables; free individual parameters retain their \(I_D\)-values.

For an \(r\)-ary term \(R\) positive in the \(r\)-ary variable \(P\), hold
the remaining interpretations fixed and define
\[
 F:\mathcal P(D^r)\to\mathcal P(D^r),\qquad
 F(X):=\llbracket R\rrbracket_{I_R[P:=X],I_D}.
\]
Then
\[
\begin{aligned}
 \llbracket\mu P[R]\rrbracket_{I_R,I_D}
   &=\bigcap\{X\subseteq D^r:F(X)\subseteq X\},\\
 \llbracket\nu P[R]\rrbracket_{I_R,I_D}
   &=\bigcup\{X\subseteq D^r:X\subseteq F(X)\}.
\end{aligned}
\]
Positivity makes \(F\) monotone. These are respectively its least and greatest
fixed points. Explicitly, put
\[
 X_0=\varnothing,\quad X_{j+1}=F(X_j),\qquad
 Y_0=D^r,\quad Y_{j+1}=F(Y_j).
\]
On finite \(D\), these chains stabilize after at most \(|D|^r\) strict changes.
The bound counts approximant changes, not the time to perform their BDD
operations.
For the least chain, induction gives \(X_j\subseteq U\) for every
pre-fixed set \(U\); at stabilization \(F(X_j)=X_j\). The descending argument
is dual. This supplies the justification for finite iteration without
assuming that an arbitrary monotone iteration from an arbitrary seed converges
to the least fixed point.

For the Boolean domain, an \(r\)-ary relation is represented by a BDD whose
placeholder variables \(d_1,\ldots,d_r\) stand for its arguments.  The source's
two mutually recursive evaluators follow the syntax of formulas and relational
terms.  Boolean connectives become BDD operations, existential quantification
becomes BDD abstraction, and applying a relation substitutes the actual
argument variables for its placeholders.  For \(\mu P[R]\), the printed
fixed-point routine starts with \(Z_0=\mathit{false}\) and repeats
\[
 Z_{i+1}
   =\operatorname{BDD}_{R}
      \bigl(R,I_R[P\mapsto Z_i]\bigr)
\]
until \(Z_{i+1}=Z_i\).  Formal monotonicity makes the approximants increasing;
finiteness makes them stabilize at the least fixed point.  Equality is cheap
because reduced BDDs under the chosen order are canonical.  An arbitrary
finite domain is handled by choosing \(m\geq1\) and a surjection
\(\phi:\{0,1\}^{m}\twoheadrightarrow D\), replacing each domain value by a
Boolean vector.  The encoding and BDD variable order can change performance,
but not the denotation
\cite[\S4 and fig.~2]{BurchEtAl1992}.

\paragraph{The two recursive BDD evaluators.}
The placeholder variables are fresh and are not individual variables of the
logical signature. Substitutions below are simultaneous; bound variables are
renamed when needed to avoid capture. A substitution denotes Boolean
composition followed by ordered reduction, not unrestricted relabelling of
vertices in a graph whose order might then be violated.
An individual variable \(z\) may itself
be read as a Boolean formula in this Boolean-domain evaluator, true exactly
when \(I_D(z)=1\).

Using the source's formula and relation evaluators, Figure~2 gives
\[
\begin{aligned}
 \operatorname{BDD}_F(z,I_R)
   &=\operatorname{BDDAtom}(z),\\
 \operatorname{BDD}_F(f\land g,I_R)
   &=\operatorname{BDDAnd}
       (\operatorname{BDD}_F(f,I_R),\operatorname{BDD}_F(g,I_R)),\\
 \operatorname{BDD}_F(\neg f,I_R)
   &=\operatorname{BDDNegate}(\operatorname{BDD}_F(f,I_R)),\\
 \operatorname{BDD}_F(\exists z[f],I_R)
   &=\operatorname{BDDExists}(z,\operatorname{BDD}_F(f,I_R)),\\
 \operatorname{BDD}_F(R(z_1,\ldots,z_r),I_R)
   &=\operatorname{BDD}_R(R,I_R)
           [d_1:=z_1,\ldots,d_r:=z_r],\\
 \operatorname{BDD}_R(P,I_R)&=I_R(P),\\
 \operatorname{BDD}_R(\lambda z_1,\ldots,z_r[f],I_R)
   &=\operatorname{BDD}_F(f,I_R)
           [z_1:=d_1,\ldots,z_r:=d_r].
\end{aligned}
\]
Here \(I_R(P)\) in the last two lines means the BDD representing the
interpreted relation, rather than its extension as a set. The \(\mu\)-case is
the \(Z_i\)-iteration already displayed. The BDD operations denote, in order,
an atomic test, conjunction, negation, and existential abstraction.
Disjunction is recovered by De Morgan's law; the relational \(\nu\)-case
uses complement duality
\[
 \nu P[R]=\neg\mu P[\neg R[P:=\neg P]],
\]
where relational negation complements an \(r\)-ary relation within \(D^r\)
\cite[\S\S3--4 and fig.~2]{BurchEtAl1992}.

For a relation term with free individual parameters, the returned BDD may
contain those parameters as well as its \(r\) placeholders. Fixing the
parameters first gives an \(r\)-ary relation; the fixed-point iteration is
performed pointwise at those parameter values. The application condition
above prevents a placeholder substitution from accidentally changing one of
these free parameters.

The Boolean encoding of a nonempty finite \(D\) also has a precise
interpretation. For the surjection \(\phi:\{0,1\}^m\to D\), represent an
\(r\)-ary relation \(Q\) by the characteristic function of
\[
 \{(b_1,\ldots,b_r):(\phi(b_1),\ldots,\phi(b_r))\in Q\}.
\]
Every bit vector has a denotation, though different vectors may denote the
same element. Surjectivity makes existential quantification exact: a domain
witness has a Boolean representative, and every Boolean witness maps to a
domain witness. Inverse images preserve complement, union, and intersection.
Induction through the finite fixed-point approximants therefore gives the
same relation after decoding, including when \(|D|\) is not a power of two.

The paper then instantiates this evaluator for a Kripke structure
\[
 M=(A,S,L,N,S_0),
\]
where \(A\) is the atomic-proposition set, \(S\) is finite,
\(L:S\to2^A\), the transition relation \(N\subseteq S\times S\) is total, and
\(S_0\subseteq S\) is the initial set.  Its primitive CTL temporal forms are
\(EXf\), \(E[f\,U\,g]\), and \(EGf\).  Reading a CTL formula as a unary
relational term gives the exact translations
\[
\begin{aligned}
 EXf
  &=\lambda s\,\exists t\,[f(t)\land N(s,t)],\\
 EGf
  &=\nu Q\,\lambda s\,[f(s)\land
       \exists t\,(Q(t)\land N(s,t))],\\
 E[f\,U\,g]
  &=\mu Q\,\lambda s\,[g(s)\lor
       (f(s)\land\exists t\,(Q(t)\land N(s,t)))].
\end{aligned}
\]
The whole structure satisfies \(f\) at its initial states exactly when
\(\forall s\,[S_0(s)\to f(s)]\) is true.  If states are encoded by current
Boolean variables \(V\) and successors by a renamed copy \(V'\), the common
predecessor operation is therefore
\[
 \operatorname{Pre}_{\exists}(X)(V)
 =\exists V'\,\bigl(N(V,V')\land X(V')\bigr),
\]
implemented by conjunction, existential abstraction, and renaming on BDDs
\cite[\S6.1]{BurchEtAl1992}.

For a concrete fixed-point trace, take total transitions
\(s_0\to s_1\), \(s_0\to s_2\), \(s_1\to s_1\), and \(s_2\to s_2\), let
\(g\) hold only at \(s_1\), and let \(f\) hold at \(s_0,s_1\).  The least
approximants for \(E[f\,U\,g]\) are
\[
 \varnothing,
 \quad\{s_1\},
 \quad\{s_0,s_1\},
 \quad\{s_0,s_1\}.
\]
Representing these sets and \(N\) by their characteristic BDDs changes only
the data structure: the second step is still conjunction with \(f\), symbolic
preimage through \(N\), and union with \(g\).  The structural evaluator is
correct because each BDD operation denotes its corresponding logical or
relational operation, and the fixed-point case is the finite monotone iteration
above.

The symbolic calculation can be made explicit with two current bits \(b,c\)
and next bits \(b',c'\). This is a handbook calculation of the preceding
example. Use the surjection
\[
 \phi(0,0)=s_0,\qquad
 \phi(0,1)=s_1,\qquad
 \phi(1,0)=\phi(1,1)=s_2.
\]
The Boolean predicates for \(f\), \(g\), and the transition relation are
\[
\begin{aligned}
 f(b,c)&=\neg b,& g(b,c)&=\neg b\land c,\\
 N(b,c,b',c')
  &=(\neg b\land\neg b'\land c')
      \lor((b\lor\neg c)\land b').
\end{aligned}
\]
The first disjunct encodes \(s_0\to s_1\) and \(s_1\to s_1\).
The second encodes \(s_0\to s_2\) and \(s_2\to s_2\), including both Boolean
representatives of \(s_2\). No bit vector is an invalid state code.

For a state predicate \(U(b,c)\), its symbolic preimage is
\[
 \operatorname{Pre}_{\exists}(U)(b,c)
 =\exists b'\exists c'\,
      (N(b,c,b',c')\land U(b',c')).
\]
Here \(U(b',c')\) is simultaneous renaming to the next-state bits.
For \(U=g\), substituting \(\neg b'\land c'\) removes the second disjunct
of \(N\), leaving \(\neg b\land\neg b'\land c'\); existential abstraction
therefore yields \(\neg b\).
The until iteration is consequently
\[
 Z_0=0,\quad Z_1=\neg b\land c,\quad
 Z_2=(\neg b\land c)\lor(\neg b\land\neg b)=\neg b.
\]
Substituting \(Z_2(b',c')=\neg b'\) again gives preimage \(\neg b\), so
\(Z_3=Z_2\). Decoding yields precisely \(\{s_0,s_1\}\).
This calculation exposes the conjunction, renaming, and quantifier
elimination behind each state-set step.

\paragraph{Positive transitive closure by squaring.}
For \(E\subseteq S\times S\), relational composition is
\[
 (Q\circ Q)(x,y)\iff
       \exists z\in S\;(Q(x,z)\land Q(z,y)).
\]
The iterative-squaring construction computes positive transitive closure:
\[
 \mu Q\,\lambda x,y\,[E(x,y)\lor
   \exists z\,(Q(x,z)\land Q(z,y))].
\]
Starting with \(Q_0=\varnothing\), put \(Q_{i+1}=E\cup(Q_i\circ Q_i)\).
For \(i\geq1\),
\[
 Q_i=\bigcup_{\ell=1}^{2^{i-1}}E^\ell.
\]
The base case is \(Q_1=E\). Composition of two positive paths of
length at most \(2^{i-1}\) produces lengths from \(2\) through \(2^i\);
every length in that interval splits into two permitted positive
lengths. The extra \(E\) supplies length \(1\), proving the induction.
Consequently the limit is \(E^+\), not the reflexive closure \(E^*\).
The latter is \(=_{S}\cup E^+\).

Let \(d_+(E)\) be the maximum shortest positive-path length over pairs
in \(E^+\), taking \(d_+(E)=0\) when \(E^+=\varnothing\). Then the
complete relation has been reached by round
\[
 1+\left\lceil\log_2\max\{1,d_+(E)\}\right\rceil.
\]
An additional equality comparison detects stabilization. A shortest
positive path between distinct states is simple; a shortest positive
return to the same state is a simple cycle. Thus \(d_+(E)\leq|S|\).
This bound includes self-return pairs, whose length is not their
zero-valued reflexive distance. The logarithmic iteration count does
not bound the sizes of the squared-relation BDDs
\cite[\S5]{BurchEtAl1992}.

\paragraph{The general iterative-squaring transformation.}
Let \(D\) be a nonempty finite domain, \(r\geq1\), and \(V=D^r\).
Fix a seed relation \(B\subseteq V\) and a transition kernel
\(K\subseteq V\times V\). In the one-domain relational calculus,
these have arities \(r\) and \(2r\), respectively. Their defining
terms may have fixed parameters, but neither may contain the
iteration variable \(Q\) free. Write
\[
 \operatorname{Post}_K(X)
   :=\{\bar y\in V:\exists\bar x\in X\;K(\bar x,\bar y)\},
 \qquad
 \mathcal F(X):=B\cup\operatorname{Post}_K(X).
\]
Thus the admitted body has the form
\[
 \lambda\bar y\,
   [B(\bar y)\lor
      \exists\bar x\,(Q(\bar x)\land K(\bar x,\bar y))].
\]
The tuples of bound variables are chosen fresh and disjoint.
The restriction that \(B,K\) do not depend on \(Q\) makes this a
fixed graph-reachability operator throughout the iteration.
It is the form used in the general transformation of
Section~5.2 \cite[\S5.2, thms.~2--3]{BurchEtAl1992}.

Compute the positive closure \(T=K^+\) by the preceding squaring
construction, now on tuple states in \(V\). Define
\[
\begin{aligned}
 M_B&=B\cup\operatorname{Post}_T(B),\\
 C_K&=\{\bar y\in V:\exists\bar x\in V\
                 [T(\bar x,\bar x)\land T(\bar x,\bar y)]\}.
\end{aligned}
\]
Then
\[
 \mu\mathcal F=M_B,\qquad
 \nu\mathcal F=M_B\cup C_K.
\]
The first set consists of the seed states and everything reachable
from them. The additional set in the greatest fixed point consists
of states reachable from a positive closed walk. A point on such
a walk is included even when the route to itself is regarded as
length zero, because its positive return already witnesses
\(T(\bar x,\bar x)\).

\emph{Proof of the least-fixed-point formula.}
The set \(M_B\) contains \(B\) and is closed under one more
\(K\)-edge. Conversely, a nonseed member of \(M_B\) has a
positive path from \(B\); deleting its final edge leaves a
predecessor in \(M_B\). These two observations give
\(\mathcal F(M_B)=M_B\).
If \(\mathcal F(X)\subseteq X\), then \(B\subseteq X\) and
\(\operatorname{Post}_K(X)\subseteq X\). Induction on path
length puts every state reachable from \(B\) in \(X\).
Thus \(M_B\) is contained in every pre-fixed set and is the
least fixed point.

\emph{Proof of the greatest-fixed-point formula.}
First,
\[
 \operatorname{Post}_K(C_K)=C_K.
\]
One inclusion follows by appending an edge to the path from
a closed walk. For the other, let a positive path from a point
\(\bar x\) on a closed walk end at \(\bar y\in C_K\).
Its last-edge predecessor is itself reachable from that walk,
and hence lies in \(C_K\). If this predecessor is \(\bar x\),
its positive self-return supplies the required membership.
Therefore every \(\bar y\in C_K\) has a predecessor in \(C_K\).
Since relational image distributes over union and \(M_B\)
is fixed,
\[
 \mathcal F(M_B\cup C_K)=M_B\cup C_K.
\]

Now let \(X\subseteq\mathcal F(X)\) and take \(\bar y_0\in X\).
If \(\bar y_0\notin B\), post-fixedness supplies
\(\bar y_1\in X\) with \(K(\bar y_1,\bar y_0)\).
Repeat this predecessor choice until a seed is encountered.
If \(\bar y_j\in B\) is reached, the selected edges give a
finite forward path from that seed to \(\bar y_0\), putting
\(\bar y_0\) in \(M_B\).
Otherwise there is an infinite predecessor sequence in the
finite set \(V\). Some \(\bar y_i=\bar y_j\), \(i<j\).
The edges from \(\bar y_j\) down to \(\bar y_i\) give a
positive closed walk, and the remaining edges lead to
\(\bar y_0\). If \(i=0\), the closed walk itself gives a
positive return to \(\bar y_0\). In every case
\(\bar y_0\in C_K\).
Thus every post-fixed set is contained in \(M_B\cup C_K\);
the fixed set already exhibited is the greatest one.
\(\square\)

Finiteness is essential to that last argument. On
\(V=\{0,1,2,\ldots\}\), take \(B=\varnothing\) and
\(K(i+1,i)\) for every \(i\). There is no cycle, so both
\(M_B\) and \(C_K\) are empty, yet
\(\operatorname{Post}_K(V)=V\) and the greatest fixed point
is all of \(V\). Infinite predecessor support need not
contain a repeated state on an infinite domain.

For a finite worked example, let
\[
 V=\{a,b,c,d,e\},\qquad B=\{a\},\qquad
 K=\{(a,b),(b,c),(d,d),(d,e)\}.
\]
The least approximants are
\[
 \varnothing,\quad\{a\},\quad\{a,b\},\quad\{a,b,c\},
 \quad\{a,b,c\}.
\]
The positive return at \(d\) and its edge to \(e\) give
\(C_K=\{d,e\}\), so the greatest fixed point is \(V\).
The state \(e\) belongs to it through an incoming edge
from the closed walk; \(e\) need not itself have an
outgoing edge. This is a forward-image fixed point.

The CTL predecessor operators fit the same form after
reversing the kernel arguments. For
\[
 \mathcal U(X)=G\cup(F\cap\operatorname{Pre}_{\exists}(X)),
\]
use seed \(G\) and kernel
\[
 K(x,y)\iff F(y)\land N(y,x).
\]
Then \(\mathcal U(X)=G\cup\operatorname{Post}_K(X)\),
so the least formula computes \(E[F\,U\,G]\).
With empty seed, the greatest formula computes
\(EG(F)\). The direction of \(K\) matters: in the
original relation \(N\), these are states that can
reach a target or continue indefinitely, not states
reached by running \(N\) forward from the target.

Monotonicity alone is not enough to justify this graph
transformation. Every operator \(B\cup\operatorname{Post}_K(X)\)
preserves binary unions. Define the universal-predecessor operator by
\[
 \operatorname{Pre}_{\forall}(X)
  :=\{s:\forall t\,(N(s,t)\Rightarrow t\in X)\}.
\]
It does not preserve binary unions: if \(s\) has exactly two distinct successors
\(a,b\), it belongs to
\(\operatorname{Pre}_{\forall}(\{a,b\})\) but not to
\(\operatorname{Pre}_{\forall}(\{a\})\cup
  \operatorname{Pre}_{\forall}(\{b\})\).
Consequently such a body cannot simply be substituted
for the admitted form; any use of duality or another
rewriting must first establish the required equality.

Once \(B,K\) are evaluated, the \(T\)-stage reaches closure
in \(O(1+\log|V|)\) rounds, followed by the displayed
Boolean, image, and diagonal operations. It computes
a \(2r\)-ary relation where the original approximants
were only \(r\)-ary. The round bound is therefore not
a running-time improvement theorem for arbitrary BDDs,
nor does it include the cost of evaluating nested terms
inside \(B\) or \(K\).

\paragraph{Applying the transformation inside a nested fixed point.}
Let \(E\subseteq V\times V\) and \(A\subseteq V\), and use
\(\operatorname{Pre}_E(X)=\{y:\exists x\in X\ E(y,x)\}\).
Consider
\[
 Z=\nu P\,
   \left[A\cap
      \mu Q\,\bigl(\operatorname{Pre}_E(P)
                     \cup\operatorname{Pre}_E(Q)\bigr)\right].
\]
For the inner iteration, \(P\) is a fixed parameter.
Its seed is \(\operatorname{Pre}_E(P)\), so the inner
least fixed point is exactly
\(\operatorname{Pre}_{E^+}(P)\): it permits a positive
number of \(E\)-steps to \(P\).
For the outer iteration, the body becomes
\(A\cap\operatorname{Pre}_{E^+}(P)\), with empty seed
and kernel \(K(x,y)\iff A(y)\land E^+(y,x)\).
It is again of the admitted form.

This nested expression has the following direct result:
\[
 Z=\{y\in A:\exists x\in A\
                       [E^+(y,x)\land E^+(x,x)]\}.
\]
To prove it, first observe that a post-fixed point of
\(P\mapsto A\cap\operatorname{Pre}_{E^+}(P)\) lets one
choose, from each of its states, a positive finite path
to another state of that set. Concatenating such paths
gives an infinite path visiting \(A\) infinitely often.
Conversely, an infinite path beginning in \(A\) and
visiting \(A\) infinitely often yields a post-fixed set
by taking the states at those visits. This proves that
the greatest fixed point denotes exactly those starting
states.

On a finite graph, some \(A\)-state then occurs twice
strictly after the starting position, giving both a
positive route from \(y\) and a positive return at \(x\).
Conversely, a route to such an \(x\), followed by
repetition of its closed walk, visits \(A\) infinitely
often. This proves the displayed formula, including
the case \(x=y\).
For example, with
\(E=\{(a,b),(b,c),(c,b)\}\) and \(A=\{a,c\}\),
the set is \(\{a,c\}\). The recurring segment passes
through \(b\notin A\); requiring every intermediate
state to remain in \(A\) would express a different
property \cite[\S5.2]{BurchEtAl1992}.

\paragraph{Fair paths as nested fixed points.}
Let \(C_1,\ldots,C_k\subseteq S\), with \(k\geq0\), be fixed state
sets. An infinite path \(\pi=s_0s_1\cdots\) is fair when
\[
 \forall j\in\{1,\ldots,k\}\ \forall r\geq0\
       \exists \ell\geq r\quad s_\ell\in C_j.
\]
Thus each constraint must be visited infinitely often, and every infinite
path is fair when the list is empty. A CTL formula supplied as a constraint
is first evaluated as a state predicate in the underlying structure.
The sets \(C_j\) are then held fixed while the fair temporal operators are
evaluated. This prevents a circular definition of the constraints.
Section~6.2 implements this repeated-visit semantics with nested least and
greatest fixed points \cite[\S6.2]{BurchEtAl1992}.

For \(F,B\subseteq S\), put
\[
 \operatorname{Reach}_F(B)
   :=\mu Y\,\bigl(B\cup(F\cap\operatorname{Pre}_{\exists}(Y))\bigr).
\]
The least approximants show that this is exactly the set of starting
states of finite paths ending in \(B\), with every state before the
endpoint in \(F\). The length-zero path is permitted; in particular
\(B\subseteq\operatorname{Reach}_F(B)\). The corresponding
positive-length constrained reachability set is
\[
 F\cap\operatorname{Pre}_{\exists}
           \bigl(\operatorname{Reach}_F(B)\bigr).
\]
The extra preimage supplies the first edge, and the intersection
requires its starting state to lie in \(F\).

Adjoin the harmless constraint \(C_0=S\), which every infinite path
satisfies, and define
\[
 \Phi_F(Z)
   :=F\cap\bigcap_{j=0}^{k}
       \operatorname{Pre}_{\exists}
          \bigl(\operatorname{Reach}_F(Z\cap C_j)\bigr),
 \qquad
 E_CG(F):=\nu Z\,\Phi_F(Z).
\]
All operations are on \(\mathcal P(S)\). Preimage and constrained
reachability are monotone in their set arguments, so \(\Phi_F\) is
monotone. At a fixed point \(Z=\Phi_F(Z)\), every state lies in \(F\)
and, for every \(j\), has a positive-length \(F\)-path to a state in
\(Z\cap C_j\). In particular, the endpoint remains in the set from
which the next visit can be arranged.

\emph{Fair-path characterization.}
The set \(E_CG(F)\) consists exactly of the states from which there is
a fair infinite path all of whose states lie in \(F\).

For one direction, write \(Z=\nu\Phi_F\) and start at \(s\in Z\).
Use its fixed-point property to choose a finite path to \(Z\cap C_0\),
then one to \(Z\cap C_1\), and so on through \(Z\cap C_k\), repeating
this schedule forever. Every chosen segment has positive length.
Its initial state and intermediate states lie in \(F\), and its
endpoint lies in \(Z\subseteq F\). Concatenation therefore produces
an infinite \(F\)-path, with a visit to every \(C_j\) during each
schedule round. It is fair.

Conversely, take a fair infinite \(F\)-path and let \(W\) be the set
of states occurring on it. From any chosen occurrence of a state
in \(W\), each \(C_j\) occurs again strictly later; \(C_0\) can be
reached at the next position. The intervening finite segment lies
in \(F\) and ends in \(W\cap C_j\). Hence \(W\subseteq\Phi_F(W)\).
The greatest fixed point contains every such post-fixed set, so the
initial state belongs to \(\nu\Phi_F\). This proves the characterization.

The added \(C_0\) also handles \(k=0\) correctly. In that case the
construction still requires indefinitely repeatable positive-length
segments and gives ordinary \(EG(F)\). An empty intersection without
this requirement would instead yield \(F\), even when an \(F\)-state
has no infinite continuation inside \(F\).
The proof works on a finite directed graph even if its transition
relation is not total: a dead end simply has no infinite-path witness.

Put \(H:=E_CG(S)\), the states with some fair continuation.
If \(F,G\) are the already evaluated operand state sets, then
\[
 E_CX(F)=\operatorname{Pre}_{\exists}(F\cap H),
 \qquad
 E_C[F\,U\,G]=\operatorname{Reach}_F(G\cap H).
\]
A fair next witness has a successor in \(F\) whose remaining suffix
is fair. A fair until witness has a finite \(F\)-prefix ending in
\(G\) and then a fair suffix. Conversely, concatenate the indicated
finite prefix with that fair suffix; finitely many added positions
cannot change infinitely many visits. These observations prove both
identities. Universal fair operators follow by duality over the same
set of fair paths. At a state with no fair paths, existential fair
assertions are false and universal ones are vacuously true.

For an explicit calculation, take transitions
\[
 a\to a,\quad a\to b,\quad b\to b,\quad c\to c,
 \qquad C_1=\{b\}.
\]
With \(F=S\), the outer iteration starts at \(S=\{a,b,c\}\).
The inner reachability set for \(C_1\) is \(\{a,b\}\), and its
preimage is again \(\{a,b\}\); the tautological constraint adds
no restriction here. The next outer set is \(\{a,b\}\), which is
fixed. Thus \(H=\{a,b\}\). The ordinary \(EG(\{a\})\) assertion
holds at \(a\) via its self-loop, but \(E_CG(\{a\})\) is empty:
that loop never visits \(b\).

Symbolically, each inner computation uses union, conjunction, and
the existing BDD preimage operation. Starting the outer iteration
at \(S\) and each inner least iteration at \(\varnothing\) computes
the displayed fixed points. There are at most \(|S|\) strict
changes in each chain. A direct nested implementation uses
\(O((k+1)(|S|+1)^2)\) preimage evaluations; this counts set
operations, not their BDD running time or intermediate graph sizes.

\paragraph{A symbolic tableau for linear temporal logic.}
Let \(f\) be an LTL formula over its finite set of atomic propositions
\(P\), with constructors \(\neg,\lor,X,U\), and use the infinite-word
semantics of Definition~\ref{def:p09-ltl}. The tableau uses one Boolean
coordinate for each member of the finite set \(\operatorname{el}(f)\):
\[
\begin{aligned}
 \operatorname{el}(p)&=\{p\},\\
 \operatorname{el}(\neg g)&=\operatorname{el}(g),\\
 \operatorname{el}(g\lor h)&=\operatorname{el}(g)\cup\operatorname{el}(h),\\
 \operatorname{el}(Xg)&=\{Xg\}\cup\operatorname{el}(g),\\
 \operatorname{el}(g\,U\,h)
   &=\{X(g\,U\,h)\}\cup\operatorname{el}(g)\cup\operatorname{el}(h).
\end{aligned}
\]
Repeated occurrences of the same syntactic formula use one coordinate.
Write \(\mathcal V_f=\{0,1\}^{\operatorname{el}(f)}\) for the
valuation space and \(v_e\) for coordinate \(e\).
The coordinate \(X(g\,U\,h)\) records a next-step obligation; the
truth of \(g\,U\,h\) at the current step is computed, not stored
as another independent bit \cite[\S7]{BurchEtAl1992}.

For each subformula \(g\) of \(f\), define the Boolean predicate
\(\alpha_g:\mathcal V_f\to\{0,1\}\) recursively:
\[
\begin{aligned}
 \alpha_p(v)&=v_p,\\
 \alpha_{\neg g}(v)&=\neg\alpha_g(v),\\
 \alpha_{g\lor h}(v)&=\alpha_g(v)\lor\alpha_h(v),\\
 \alpha_{Xg}(v)&=v_{Xg},\\
 \alpha_{g\,U\,h}(v)
   &=\alpha_h(v)\lor
           \bigl(\alpha_g(v)\land v_{X(g\,U\,h)}\bigr).
\end{aligned}
\]
The last clause recurses only into \(g,h\); the remaining expression
is a coordinate lookup. Thus the recursion is well founded even
though an until formula occurs inside its next-obligation coordinate.
For a generated coordinate \(Xg\), the predicate \(\alpha_g\)
is defined by the same clauses, including when \(g\) is the until
formula that generated that coordinate.

The initial predicate and transition relation are
\[
 I_f(v):=\alpha_f(v),\qquad
 R_f(v,w):=
   \bigwedge_{Xg\in\operatorname{el}(f)}
        \bigl(v_{Xg}\leftrightarrow\alpha_g(w)\bigr).
\]
These are finite Boolean expressions on current and next coordinates,
so their BDDs can be constructed by the operations already given.
For every until subformula \(u=g\,U\,h\), add the fairness set
\[
 C_u:=\{v\in\mathcal V_f:
              \neg\alpha_u(v)\lor\alpha_h(v)\}.
\]
It requires infinitely many positions at which that until obligation
is either absent or fulfilled. It does not require its right operand
to recur after the obligation has ceased.

\emph{Truth lemma.}
If \(v^0v^1\cdots\) is an \(R_f\)-path fair for all \(C_u\), define
the word \(\sigma\) by
\(\sigma(i)=\{p\in P:v^i_p=1\}\). Then, for every subformula \(g\)
and every \(i\geq0\),
\[
 \alpha_g(v^i)=1\quad\Longleftrightarrow\quad\sigma,i\models g.
\]
The atomic and Boolean cases follow from their definitions. For \(Xg\),
the transition equation gives
\(\alpha_{Xg}(v^i)=v^i_{Xg}=\alpha_g(v^{i+1})\), so induction
gives the next clause.
For \(u=g\,U\,h\), put
\[
 w_i=\alpha_u(v^i),\qquad
 b_i=\alpha_h(v^i),\qquad a_i=\alpha_g(v^i).
\]
The predicate and transition equations give
\[
 w_i=b_i\lor(a_i\land w_{i+1}).
\]
If \(w_i=1\) and no \(b_j\) with \(j\geq i\) were true, this
recurrence would force \(w_j=a_j=1\) for every \(j\geq i\).
The path would then avoid \(C_u\) forever, contrary to fairness.
Let \(j\geq i\) be the first position with \(b_j=1\).
The recurrence forces \(a_\ell=1\) for \(i\leq\ell<j\), and
induction on the operands gives a semantic until witness.
Conversely, any semantic until witness gives \(b_j=1\) and those
intervening \(a_\ell=1\). Starting with \(w_j=1\) and using the
recurrence backwards gives \(w_i=1\). This proves the until case
and the lemma.

For the converse construction, start with any word \(\sigma\).
Set \(v^i_p\) to its atomic truth value and set every next-coordinate
\(v^i_{Xg}\) to the truth value of \(\sigma,i+1\models g\).
The usual until unfolding and structural induction then give
\(\alpha_g(v^i)=1\) exactly when \(\sigma,i\models g\).
Consequently every \(R_f(v^i,v^{i+1})\) holds. Each \(C_u\)
occurs infinitely often: otherwise, from some point onward \(u\)
would be true while its right operand remained false forever,
contradicting the semantics of until.

It follows in both directions that
\[
 f\text{ is satisfiable}
 \quad\Longleftrightarrow\quad
 \exists v\in\mathcal V_f\
     \bigl(I_f(v)\land v\in E_CG(\mathcal V_f)\bigr).
\]
Here \(E_CG\) is evaluated on \(R_f\), with the displayed until
fairness sets and the tautological constraint used above. In
particular, satisfiability needs \emph{some} initial tableau state
with a fair continuation, not every valuation satisfying \(I_f\).
The relation \(R_f\) need not be total. For example, independent
coordinates for \(Xp\) and \(X\neg p\) may both be set to \(1\),
leaving no possible successor. The infinite-path fixed point
correctly discards that state; adding a self-loop at a dead end
would destroy the transition equations and the truth lemma.

For \(f=p\,U\,q\), the coordinates are \(p,q,X(p\,U\,q)\).
Writing the last bit as \(t\), the construction is
\[
 I_f(p,q,t)=q\lor(p\land t),\qquad
 R_f(v,w)\iff t=q'\lor(p'\land t'),\qquad
 C_f(v)\iff\neg(q\lor(p\land t))\lor q.
\]
The constant path \((1,0,1)^\omega\) satisfies both the initial
and transition predicates, but not fairness: it postpones \(q\)
forever. In contrast,
\[
 (1,0,1)\,(0,1,0)\,(0,0,0)^\omega
\]
is an initial fair path. The until obligation is fulfilled at
position \(1\); all later states satisfy \(C_f\) because the
obligation is false. The example distinguishes fulfilling an
eventuality from demanding its right operand infinitely often.

There are \(|\operatorname{el}(f)|=O(|f|)\) Boolean coordinates
and at most \(2^{|\operatorname{el}(f)|}\) tableau states. A BDD
represents their sets and transition constraints without requiring
an explicit list of those states; no small-BDD bound follows
solely from the linear number of coordinates.

\paragraph{Strong and weak observational equivalence.}
Consider two finite labelled transition systems with the same finite
action alphabet:
\[
 \mathcal M_S=(S,s_*,\Sigma,\Delta_S),\qquad
 \mathcal M_T=(T,t_*,\Sigma,\Delta_T),
\]
where \(\Delta_S\subseteq S\times\Sigma\times S\) and
\(\Delta_T\subseteq T\times\Sigma\times T\). Transitions need not
be total. For \(Z\subseteq S\times T\), define
\[
\begin{split}
 \mathcal B(Z)=\{(s,t)\in S\times T:\;&
   \forall a\in\Sigma\ \forall s'\in S\
     [\Delta_S(s,a,s')\Rightarrow
       \exists t'\in T\,(\Delta_T(t,a,t')\land Z(s',t'))],\\
 &\forall a\in\Sigma\ \forall t'\in T\
     [\Delta_T(t,a,t')\Rightarrow
       \exists s'\in S\,(\Delta_S(s,a,s')\land Z(s',t'))]\}.
\end{split}
\]
A strong bisimulation is precisely a relation \(Z\) with
\(Z\subseteq\mathcal B(Z)\). The operator is monotone, since enlarging
\(Z\) preserves all matching-move witnesses. Therefore
\[
 {\sim}\ :=\nu Z\,\mathcal B(Z)
\]
is itself a strong bisimulation and contains every strong bisimulation.
The two pointed systems are strongly equivalent exactly when
\((s_*,t_*)\in{\sim}\). This proves the greatest-fixed-point
characterization, rather than merely giving a recurrence.
Descending iteration from \(S\times T\) takes at most \(|S||T|\)
strict changes. The universal and existential clauses are relational
Boolean operations, so the existing symbolic evaluator applies
\cite[\S8.1]{BurchEtAl1992}.

For weak equivalence, designate an internal action \(\tau\in\Sigma\).
Its zero-or-more-step relation in the first system is
\[
 H_S=\mu Y\,\lambda s,t\,
      [s=t\lor\exists u\in S\,
          (\Delta_S(s,\tau,u)\land Y(u,t))],
\]
and define \(H_T\) analogously. Least iteration gives exactly the
reflexive transitive closure of the internal transitions.
Define saturated transitions by
\[
 \widehat\Delta_S(s,a,t)\iff
 \begin{cases}
  H_S(s,t),&a=\tau,\\
  \exists u,v\in S\
   [H_S(s,u)\land\Delta_S(u,a,v)\land H_S(v,t)],&a\ne\tau.
 \end{cases}
\]
Use the corresponding definition in \(T\). A visible action is
surrounded by any finite number of internal steps; an internal
action may be matched by no step at all.
The weak bisimilarity relation is
\[
 {\approx}\ :=\nu Z\,\widehat{\mathcal B}(Z),
\]
where \(\widehat{\mathcal B}\) is the previous two-sided matching
operator with both transition relations replaced by their saturated
versions \cite[\S8.2]{BurchEtAl1992}.

To justify this characterization, define a weak bisimulation directly:
each original transition on either side must have a matching saturated
transition on the other, ending in a related pair. A bisimulation for
the saturated systems has this property because every original
transition is saturated. Conversely, suppose \(Z\) has the direct
weak-matching property. Match a finite string of internal moves
one move at a time. The successive endpoints remain \(Z\)-related,
and concatenating the matching internal paths gives an \(H\)-path.
For a visible saturated move, apply this argument to its internal
prefix, match the single visible move, and apply it to the internal
suffix. The concatenated witness is the required saturated move.
Zero internal moves are matched by reflexivity. Thus
\(Z\subseteq\widehat{\mathcal B}(Z)\), proving that the two
characterizations have the same greatest relation.

This is divergence-insensitive weak equivalence: a state with only an
internal self-loop and a state with no transitions are weakly
equivalent. Their saturated internal relations both permit stuttering,
and neither has a visible move. They are not strongly equivalent,
because the original internal move has no strong match.
Replacing \(H_S\) by a positive internal closure in the \(\tau\)-case
would incorrectly exclude the zero-step match.

To use the one-domain relational calculus, take disjoint copies of
\(S,T,\Sigma\) inside \(D=S\sqcup T\sqcup\Sigma\). Interpret the
transition predicates on their displayed sorts and express each
restricted quantifier with the corresponding unary membership
predicate. Equality supplies the diagonal in \(H_S,H_T\).
This makes the state-pair and action arguments explicit before
Boolean encoding; the fixed point for strong or weak equivalence
is over the state-pair relation.

\paragraph{Language containment with a deterministic specification.}
Let
\[
 \mathcal A=(S,s_*,\Sigma,\Delta_A,F_A),\qquad
 \mathcal D=(T,t_*,\Sigma,\delta_D,F_D)
\]
be finite B\"uchi automata. The first may be nondeterministic;
the second is deterministic and complete:
\(\delta_D:T\times\Sigma\to T\) is a function defined on every
state-letter pair. A run starts at the designated initial state
and is accepting exactly when it visits its accepting set
infinitely often. Use product states \(U=S\times T\), initial
state \(u_*=(s_*,t_*)\), and edges
\[
 R_\times((s,t),(s',t'))
   \iff\exists a\in\Sigma\
       [\Delta_A(s,a,s')\land t'=\delta_D(t,a)].
\]
The two acceptance predicates on the product are
\[
 P_A=F_A\times T,\qquad P_D=S\times F_D.
\]
Then
\[
 L(\mathcal A)\subseteq L(\mathcal D)
 \quad\Longleftrightarrow\quad
 \text{every infinite product path from }u_*
 \text{ satisfies }GF\,P_A\Rightarrow GF\,P_D.
\]
Indeed, a word accepted by \(\mathcal A\) has an accepting run.
Completeness supplies the run of \(\mathcal D\) on that same
word, and determinism makes it the only run. Pairing them shows
that a word outside \(L(\mathcal D)\) gives a product path visiting
\(P_A\) infinitely often but \(P_D\) only finitely often.
Conversely, choose a witnessing letter for each edge of such
a product path. Its first projection is accepting, and its
second projection is the unique run on the resulting word and
is rejecting. This proves the equivalence
\cite[\S9]{BurchEtAl1992}.

Both restrictions on the specification automaton matter. A missing
transition can be completed by adding a nonaccepting sink before
forming the product. Without determinism, a rejecting specification
run need not mean that the word is rejected. For example, over
\(\{a\}\), an automaton with an accepting initial state \(u\),
a nonaccepting state \(v\), and edges
\(u\to u,u\to v,v\to v\) accepts \(a^\omega\) but also has
a rejecting run on it. Universal quantification over paired runs
would reject containment of \(\{a^\omega\}\) in its language.

The containment test has two symbolic forms. First, use the fair
greatest-fixed-point construction with the single fairness set
\(P_A\), and put
\[
 W:=E_{\{P_A\}}G(U\setminus P_D).
\]
It denotes states with an infinite continuation that avoids
\(P_D\) entirely and visits \(P_A\) infinitely often. A path
visiting \(P_D\) only finitely often has such a suffix; adding
an arbitrary finite prefix does not affect either recurrence
condition. Hence
\[
 L(\mathcal A)\subseteq L(\mathcal D)
 \quad\Longleftrightarrow\quad
 u_*\notin\operatorname{Reach}_U(W),
\]
where preimages in \(\operatorname{Reach}\) now use \(R_\times\).
This formulation also works when the product has dead ends,
since it asks for an infinite witness explicitly.

Alternatively, restrict the product edges to
\[
 T_{\neg D}
   :=\{(u,v)\in R_\times:
                   u\notin P_D\text{ and }v\notin P_D\},
\]
compute its positive closure \(T_{\neg D}^{+}\) by squaring, and set
\[
 B:=\{u\in P_A:T_{\neg D}^{+}(u,u)\}.
\]
Then the same containment criterion is
\[
 u_*\notin\operatorname{Reach}_U(B).
\]
To prove it, suppose an infinite violating path exists. After its
last \(P_D\)-visit it stays outside \(P_D\). Finiteness and
infinitely many \(P_A\)-visits give a state of \(P_A\) occurring
at two distinct positions in that suffix. The intervening
positive closed walk witnesses \(T_{\neg D}^{+}(u,u)\), and
the original prefix reaches \(u\in B\).
Conversely, reach a state \(u\in B\) and repeat its positive
closed walk. It avoids \(P_D\) forever and visits \(P_A\)
on every return, so it violates containment.

Both the accepting-state conjunct \(u\in P_A\) and the positive
length condition are essential. An arbitrary rejecting cycle
need not correspond to a word accepted by \(\mathcal A\); a
length-zero return supplies no infinite path.
The nested unary fixed point and the binary-closure method
therefore decide the same condition, with different intermediate
relations and no uniform guarantee that the method with fewer
iterations uses smaller BDDs.

\paragraph{The pipeline experiment.}
The empirical presentation is a three-stage synchronous pipeline: operands are
read from a four-register file, an ALU operation is performed, and the result
is written back after three clocks.  The instruction supplies source and
destination register addresses; a \(\mathsf{STALL}\) input propagates a
no-operation; and a bypass path permits immediate reuse of an ALU result.  For
an exclusive-or instruction with source fields \(a,b\) and destination field
\(c\), its latency-aware path assertion is
\[
 (X^3R)_c=(X^2R)_a\mathbin{\oplus}(X^2R)_b.
\]
Here \(R\) is a register-file value, not one Boolean proposition;
\(X^dR\) reads that value \(d\) positions along a path. The conversion
to CTL is bitwise. For word width \(w\geq1\), let \(r_{j,\ell}\)
be the atomic predicate that bit \(\ell\) of register \(j\) is \(1\),
where \(0\leq j<4\) and \(0\leq\ell<w\).
First check, for \(d=2,3\),
\[
 AG\bigwedge_{j=0}^{3}\bigwedge_{\ell=0}^{w-1}
     \bigl((EX)^d r_{j,\ell}
              \leftrightarrow(AX)^d r_{j,\ell}\bigr).
\]
The notation \((EX)^d\) means \(d\) successive existential-next
operators, and likewise for \((AX)^d\). Totality guarantees a
length-\(d\) continuation. At any reachable state, the displayed
equivalence for a bit says either that every such continuation
sets it to \(1\), or that none does. Checking all bits therefore
proves that the complete future register file is uniquely
determined at that horizon.

Define the word-valued state abbreviation \(V_j^{(d)}\) by taking
its bit \(\ell\) to be the truth value of \((EX)^d r_{j,\ell}\).
After the two checks, it equals register \(j\)'s value at horizon
\(d\) on every path from the current state. Let \(a,b,c\in
\{0,1,2,3\}\) be the current instruction's source and destination
addresses. Register selection such as \(V_a^{(2)}\) is a finite
multiplexer on the current address bits; word equality and
exclusive-or are also Boolean abbreviations on their bits.
The exclusive-or pipeline requirement can now be written in CTL:
\[
\begin{split}
 AG\biggl[&
  \neg\mathsf{STALL}\ \Rightarrow\
   \biggl(V_c^{(3)}=V_a^{(2)}\mathbin{\oplus}V_b^{(2)}
    \ \land\
    \bigwedge_{j=0}^{3}
       (c\ne j\Rightarrow V_j^{(3)}=V_j^{(2)})\biggr)
 \biggr]\\
 {}\land\ AG\biggl[&
  \mathsf{STALL}\ \Rightarrow\
     \bigwedge_{j=0}^{3}(V_j^{(3)}=V_j^{(2)})
 \biggr].
\end{split}
\]
The destination update, preservation of nondestination registers,
and stalled no-operation are thus explicit formulas. Under the
proved uniqueness conditions, each displayed state-level clause
is equivalent to its corresponding path-level clause. This does
not permit an unrestricted replacement of path operators by
existential state operators
\cite[\S10, pp.~164--166]{BurchEtAl1992}.

The measured verifier operates directly on CTL, avoiding the overhead
of translating each property through the general relational calculus.
The reported scale is equally specific. Register widths range from one to
twelve bits, with exclusive-or, addition, or both as ALU choices.  The
twelve-bit pipeline has approximately \(1.5\times10^{29}\) reachable states.
For the eight-bit exclusive-or pipeline, the transition-relation BDD has about
42,000 nodes and verification took about 22 minutes on a Sun~3/60.  Grouping
BDD variables by bit position leaves only a fixed amount of carry/control
information between groups, producing the observed linear BDD growth in this
family.  The source expressly says that this is heuristic and does not improve
the asymptotic complexity of model checking; multiplication, barrel shifting,
or another operation carrying unbounded information between bit positions can
make the diagrams unmanageable.  Symbolic representation changes the practical
representation of these systems, not the worst-case state-explosion boundary
\cite[\S\S10--11]{BurchEtAl1992}.

\subsubsection{McMillan's 1993 book and the earlier thesis}
\label{sec:p09-n-c-017-mcmillan}

McMillan's book is the registered full SMV, symbolic-algorithm, hardware/cache,
induction, equivalence, and partial-order presentation. Its chapter ranges and
scope are established, but the revised 1993 body needed for exact SMV syntax,
partitioned transitions, algorithms, examples, and proofs is not part of the
available source basis. The PDF linked from the author's thesis/book page is
not that revision: its own title page identifies it as the May 1992 Carnegie
Mellon doctoral thesis CMU-CS-92-131 \cite{McMillan1992Thesis}. That thesis and
the Burch et al. article are related to, but are not the same source state as,
the registered book \cite{McMillan1993}. The other presentations in this
section do not supply that book presentation.

\subsubsection{Biere--Cimatti--Clarke--Zhu bounded semantics}
\label{sec:p09-n-c-017-biere}

The source for the definitions below is the complete CMU technical report
CMU-CS-99-101, dated 4 January 1999 and marked ``Submitted for TACAS'99.''  The
published LNCS article has the same title and authors but a different container
and pagination; the definitions and theorem numbers below are locators in the
complete author-hosted report, not unverified claims about the typeset LNCS
pages \cite{BiereCimattiClarkeZhu1999}.

Definition~1 uses a finite Kripke structure
\[
 M=(S,I,T,L),
\]
where \(I\subseteq S\), \(T\subseteq S\times S\) is total, and
\(L:S\to2^{AP}\).  For the encoding, \(S=\{0,1\}^n\); propositional formulas
\(I(s)\), \(T(s,t)\), and \(p(s)\) recognize initial states, transitions, and
atomic labels.  LTL is put in negation normal form with atomic negation,
conjunction, disjunction, \(X,F,G,U\), and release \(R\), and retains its
ordinary semantics on infinite paths.  Existential validity
\(M\models_E h\) means that some path beginning in \(I\) satisfies \(h\);
universal validity quantifies over all such paths
\cite[defs.~1--3]{BiereCimattiClarkeZhu1999}.

Fix \(k\in\mathbb N\) and \(0\leq\ell\leq k\). A finite prefix
\(s_0\ldots s_k\) with
\(T(s_k,s_\ell)\) represents the infinite path
\[
 s_0\ldots s_{\ell-1}(s_\ell\ldots s_k)^\omega;
\]
the source calls this a \((k,\ell)\)-loop.  On such a lasso, bounded truth is
ordinary infinite-path truth. For an infinite path that is not a \(k\)-loop,
Definition~6 uses only the observed horizon. At a position \(0\leq i\leq k\),
\(Fh\) needs a witness
between \(i\) and \(k\); \(Xh\) requires \(i<k\) and truth at \(i+1\); and
\(h_1\mathbin{U}h_2\) needs a position \(j\leq k\) where \(h_2\) holds with
\(h_1\) at every earlier position.  By contrast, \(Gh\) is false without a
loop, because the prefix cannot certify the unobserved future.  Release also
needs a finite witness rather than allowing its right operand to continue
forever.  Consequently, the usual \(G/F\) and \(R/U\) dualities do not hold in
this no-loop bounded semantics.  This is deliberate underapproximation, not a
finite-trace redefinition of LTL
\cite[defs.~4--6]{BiereCimattiClarkeZhu1999}.

Lemma~7 proves soundness:
\[
 \pi\models_k h\quad\Longrightarrow\quad\pi\models h.
\]
The loop case is immediate from the definition.  In the no-loop case the
source proves the stronger position-indexed assertion by structural induction;
the release case uses its finite witness to obtain the ordinary disjunction
``the right operand forever, or the left operand eventually while the right
operand holds through that point.''  Lemma~8 proves the converse existence
claim. Standard LTL model checking forms the product of \(M\) with a
finite automaton for \(h\). An accepting product run can be chosen ultimately
periodic. For an ordinary B\"uchi automaton with state set \(Q\), the bound
is \(k\leq |S||Q|-1\), as proved below; the actual \(Q\) includes every
state introduced by the acceptance construction. Projecting the accepting
lasso back to \(M\) gives bounded truth.
Theorem~9 combines the two lemmas:
\[
 M\models_E h
 \quad\Longleftrightarrow\quad
 \exists k\in\mathbb N\;M\models_E^k h.
\]
Thus bounded truth at one \(k\) is sound; completeness quantifies over all
bounds \cite[lems.~7--8 and thm.~9]{BiereCimattiClarkeZhu1999}.

The propositional encoding first unfolds the model:
\[
 \llbracket M\rrbracket_k
 =I(s_0)\land\bigwedge_{i=0}^{k-1}T(s_i,s_{i+1}).
\]
For each possible loop start put
\[
 {}_\ell L_k:=T(s_k,s_\ell),
 \qquad
 L_k:=\bigvee_{\ell=0}^{k}{}_\ell L_k.
\]
Definitions~11 and 13 supply different recursive translations for a prefix
without a selected loop and for a selected \((k,\ell)\)-loop. In both, the
formula is in negation normal form, \(0\leq i\leq k\), an empty conjunction
means \(\mathsf{true}\), and an empty disjunction means \(\mathsf{false}\).

\paragraph{The complete no-loop translation.}
The notation \(\llbracket f\rrbracket_k^i\) denotes a propositional formula
over the state-vector copies \(s_0,\ldots,s_k\). Definition~11 is
\[
\begin{aligned}
 \llbracket p\rrbracket_k^i&=p(s_i),&
 \llbracket\neg p\rrbracket_k^i&=\neg p(s_i),\\
 \llbracket f\land g\rrbracket_k^i
   &=\llbracket f\rrbracket_k^i\land\llbracket g\rrbracket_k^i,&
 \llbracket f\lor g\rrbracket_k^i
   &=\llbracket f\rrbracket_k^i\lor\llbracket g\rrbracket_k^i,\\
 \llbracket Gf\rrbracket_k^i&=\mathsf{false},&
 \llbracket Ff\rrbracket_k^i&=\bigvee_{j=i}^{k}\llbracket f\rrbracket_k^j,\\
 \llbracket Xf\rrbracket_k^i
   &=\begin{cases}
        \llbracket f\rrbracket_k^{i+1},&i<k,\\
        \mathsf{false},&i=k.
      \end{cases}
\end{aligned}
\]
The two binary temporal clauses are
\[
\begin{aligned}
 \llbracket f\,U\,g\rrbracket_k^i
 &=\bigvee_{j=i}^{k}
      \left(\llbracket g\rrbracket_k^j\land
            \bigwedge_{q=i}^{j-1}\llbracket f\rrbracket_k^q\right),\\
 \llbracket f\,R\,g\rrbracket_k^i
 &=\bigvee_{j=i}^{k}
      \left(\llbracket f\rrbracket_k^j\land
            \bigwedge_{q=i}^{j}\llbracket g\rrbracket_k^q\right).
\end{aligned}
\]
The release witness includes \(g\) at its endpoint; until does not require
\(f\) at its endpoint. At \(i=k\), until reduces to \(g(s_k)\) and release
reduces to \(f(s_k)\land g(s_k)\) when \(f,g\) are atoms. These are the finite
witness conditions of Definition~6, with recursive bounded truth for
non-atomic operands \cite[defs.~6, 11]{BiereCimattiClarkeZhu1999}.

\paragraph{The complete selected-loop translation.}
For fixed \(\ell\), the source's successor convention is
\[
 \operatorname{succ}(i)=
 \begin{cases}i+1,&i<k,\\ \ell,&i=k.\end{cases}
\]
The formula \({}_\ell\llbracket f\rrbracket_k^i\) uses this successor and
the selected infinite lasso. Its atomic and Boolean clauses are
\[
\begin{aligned}
 {}_\ell\llbracket p\rrbracket_k^i&=p(s_i),&
 {}_\ell\llbracket\neg p\rrbracket_k^i&=\neg p(s_i),\\
 {}_\ell\llbracket f\land g\rrbracket_k^i
  &={}_\ell\llbracket f\rrbracket_k^i
     \land{}_\ell\llbracket g\rrbracket_k^i,&
 {}_\ell\llbracket f\lor g\rrbracket_k^i
  &={}_\ell\llbracket f\rrbracket_k^i
     \lor{}_\ell\llbracket g\rrbracket_k^i.
\end{aligned}
\]
Its unary temporal clauses are
\[
\begin{aligned}
 {}_\ell\llbracket Xf\rrbracket_k^i
   &={}_\ell\llbracket f\rrbracket_k^{\operatorname{succ}(i)},\\
 {}_\ell\llbracket Ff\rrbracket_k^i
   &=\bigvee_{j=\min(i,\ell)}^{k}{}_\ell\llbracket f\rrbracket_k^j,\\
 {}_\ell\llbracket Gf\rrbracket_k^i
   &=\bigwedge_{j=\min(i,\ell)}^{k}{}_\ell\llbracket f\rrbracket_k^j.
\end{aligned}
\]
The binary clauses retain both a witness before wrapping and a witness after
wrapping:
\[
\begin{split}
 {}_\ell\llbracket f\,U\,g\rrbracket_k^i
 ={}&
 \bigvee_{j=i}^{k}
   \left({}_\ell\llbracket g\rrbracket_k^j\land
         \bigwedge_{q=i}^{j-1}{}_\ell\llbracket f\rrbracket_k^q\right)\\
 &{}\lor
 \bigvee_{j=\ell}^{i-1}
   \left({}_\ell\llbracket g\rrbracket_k^j\land
         \bigwedge_{q=i}^{k}{}_\ell\llbracket f\rrbracket_k^q\land
         \bigwedge_{q=\ell}^{j-1}{}_\ell\llbracket f\rrbracket_k^q\right).
\end{split}
\]
\[
\begin{split}
 {}_\ell\llbracket f\,R\,g\rrbracket_k^i
 ={}&
 \bigwedge_{j=\min(i,\ell)}^{k}{}_\ell\llbracket g\rrbracket_k^j\\
 &{}\lor
 \bigvee_{j=i}^{k}
   \left({}_\ell\llbracket f\rrbracket_k^j\land
         \bigwedge_{q=i}^{j}{}_\ell\llbracket g\rrbracket_k^q\right)\\
 &{}\lor
 \bigvee_{j=\ell}^{i-1}
   \left({}_\ell\llbracket f\rrbracket_k^j\land
         \bigwedge_{q=i}^{k}{}_\ell\llbracket g\rrbracket_k^q\land
         \bigwedge_{q=\ell}^{j}{}_\ell\llbracket g\rrbracket_k^q\right).
\end{split}
\]
Unlike no-loop release, this clause also permits \(g\) forever.
When \(i<\ell\), the wrapped-witness disjunctions are empty: every reachable
loop position already lies in the first traversal from \(i\) to \(k\).
When \(i\geq\ell\), they cover the previously unvisited portion
\(\ell,\ldots,i-1\) of that loop \cite[defs.~12--13]{BiereCimattiClarkeZhu1999}.

\paragraph{Why a finite loop traversal suffices.}
For any fixed loop position, its successive visits have identical infinite
labelled suffixes. Every LTL subformula therefore has the same truth value at
those visits. If an until witness occurs only after a complete extra loop,
move it to the first visit to that position. Its right operand still holds,
and the new required left-operand prefix is shorter. Thus the two displayed
witness ranges are exhaustive. For release, either the right operand holds
at every reachable position, or a finite left-operand witness is required
before any failure of the right operand. Moving such a witness to its first
visit again preserves every required position, now including the endpoint.

This reasoning also justifies arbitrary nesting of temporal operators:
structural induction uses the already established meanings of the proper
subformulas at each visited position. The \(X\)-clause changes position but
reduces formula size, so it creates no circular definition.
For \(k=\ell=i=0\), the next position is \(0\), until reduces to its right
operand, and release reduces to its right operand; an infinite constant-state
lasso is being evaluated. Without that self-loop the no-loop release instead
requires both operands at position \(0\).

A wrapped witness illustrates why the second disjunction is needed. Take
\(k=2\), \(\ell=1\), and evaluate at \(i=2\), with \(p\) true only at
\(s_2\) and \(q\) true only at \(s_1\). The selected suffix is
\(s_2s_1s_2s_1\ldots\). In
\({}_1\llbracket p\,U\,q\rrbracket_2^2\), the first disjunction is false,
because \(q(s_2)\) is false. The wrapped term \(j=1\) is
\[
 q(s_1)\land p(s_2)\land\mathsf{true},
\]
and is true. In contrast,
\({}_1\llbracket q\,R\,p\rrbracket_2^2\) is false: the perpetual-\(p\) clause
fails at \(s_1\), and the wrapped \(q\)-witness also requires \(p(s_1)\).
Release cannot use a witness at the first position where its right operand
already fails.

Definition~15 combines the two translations:
\[
 \llbracket M,h\rrbracket_k
 =\llbracket M\rrbracket_k\land
 \left(
   (\neg L_k\land\llbracket h\rrbracket^0_k)
   \lor
   \bigvee_{\ell=0}^{k}
    ({}_\ell L_k\land{}_\ell\llbracket h\rrbracket^0_k)
 \right).
\]
For example, take \(S=\{0,1\}\), \(I=\{0\}\), transitions
\(0\to1\) and \(1\to1\), let \(p\) hold only at \(0\), and let \(q\) hold
only at \(1\).  For \(h=p\mathbin{U}q\):
\begin{itemize}
\item at \(k=0\), \(s_0=0\), the only loop selector is
  \(T(0,0)=\mathsf{false}\), and the no-loop until translation has no
  \(q\)-witness, so the encoding is unsatisfiable;
\item at \(k=1\), the unfolding forces \((s_0,s_1)=(0,1)\), and
  \({}_1L_1=T(1,1)\) closes the lasso; \(p\) holds before the position where
  \(q\) holds, so the encoding is satisfiable;
\item at \(k=2\), the unfolding is \((0,1,1)\), both
  \({}_1L_2\) and \({}_2L_2\) are true, and either chosen lasso satisfies the
  same until obligation.
\end{itemize}
This displays separately the transition unfolding, loop choice, and temporal
constraint; none is hidden inside a BDD operation.

The encoding has the exact correctness property
\[
 \llbracket M,h\rrbracket_k\text{ is satisfiable}
 \quad\Longleftrightarrow\quad M\models_E^k h.
\]
For the forward direction, a satisfying assignment supplies an initial state
and \(k\) transitions. If it satisfies a loop disjunct, the selected back edge
closes an infinite lasso. Induction through the selected-loop translation
shows that this lasso satisfies \(h\). If it satisfies the no-loop disjunct,
induction through the finite translation gives bounded truth of \(h\);
totality extends the prefix to an infinite execution.

Conversely, take a bounded witness. If it is a \(k\)-loop, assign its prefix
to the state vectors and select its closing edge. The loop-translation
induction makes the corresponding disjunct true. If its bounded truth uses
the finite semantics and its prefix has no closing edge, use the no-loop
disjunct. If that prefix does have a closing edge, select any such edge
instead: every no-loop witness uses only positions through \(k\), and
structural induction on the negation-normal formula shows that its truth
survives every infinite continuation of that prefix. The selected lasso
therefore satisfies \(h\), and its loop disjunct is true. This covers both
possibilities for a finite witness prefix and proves the converse
\cite[thm.~16]{BiereCimattiClarkeZhu1999}.

For a universal LTL property \(\varphi\), put its negation into negation
normal form \(h\). Correctness and the bounded-witness existence theorem give
\[
 M\models_A\varphi
 \quad\Longleftrightarrow\quad
 \llbracket M,h\rrbracket_k\text{ is unsatisfiable for every }k\in\mathbb N.
\]
The finite criteria below replace this infinite family by justified stopping
bounds. With common subformulas
shared, the report gives a translation polynomial in \(\lvert h\rvert\),
quadratic in \(k\), and linear in the sizes of the propositional encodings of
\(I,T\), and the atomic labels
\cite[defs.~10--15, thm.~16, and cor.~17]{BiereCimattiClarkeZhu1999}.

\paragraph{Searching successive horizons.}
\label{sec:p09-n-c-017-biere-horizon-search}
For a universal property \(\varphi\), let \(h\) be the
negation-normal form of \(\neg\varphi\). The increasing-horizon
procedure has the following three steps.
\begin{enumerate}
\item Starting at \(k=0\), construct \(\llbracket M,h\rrbracket_k\)
directly from the state encoding and temporal clauses, and ask a
SAT solver whether it is satisfiable.
\item If it is satisfiable, read the state vectors and a true
disjunct from the assignment. A loop disjunct gives a concrete
infinite lasso. A no-loop disjunct gives a finite witness whose
truth survives an infinite continuation by totality and bounded
soundness. Return this counterexample to \(\varphi\).
\item If it is unsatisfiable, stop with a proof of \(\varphi\)
only when a justified completeness bound has been reached.
Otherwise increment \(k\) and repeat.
\end{enumerate}
A solver timeout or resource limit is an unknown result, not
unsatisfiability. If every smaller horizon was checked and found
unsatisfiable, the first successful \(k\) is the least satisfiable
bounded horizon: encoding correctness excludes a bounded witness
at any smaller value. This does not separately minimize the
period or the number of distinct states. Without a completeness
bound, an unsuccessful search need not terminate
\cite[\S1 and opening of \S5]{BiereCimattiClarkeZhu1999}.

The fixed-\(k\) translation itself requires no tableau or
automaton: it is the direct propositional construction already
given. The automaton product used in the completeness argument
supplies a sufficient stopping bound; it is not a prerequisite
for bounded counterexample search. A known bound can also justify
a single-horizon check when the lasso-padding argument below
applies.

\paragraph{The three-bit shift-register calculation.}
\label{sec:p09-n-c-017-biere-shift-register}
The report's initial example has state
\(x=(x[0],x[1],x[2])\in\{0,1\}^3\), arbitrary initial contents,
and transition predicate
\[
\begin{aligned}
 I(x)&=\mathsf{true},\\
 T(x,x')&\iff
   x'[0]=x[1]\ \land\ x'[1]=x[2]\ \land\ x'[2]=1.
\end{aligned}
\]
The intended zero-filling shift would use \(x'[2]=0\) and
reach the zero state after three shifts. Instead, the displayed
device evolves as
\[
 (a,b,c)\longrightarrow(b,c,1)
 \longrightarrow(c,1,1)\longrightarrow(1,1,1),
\]
and \((1,1,1)\) is a fixed point. Define the nonzero predicate
\[
 \mathsf{NZ}(x):=x[0]\lor x[1]\lor x[2].
\]
The universal eventual-zero claim is
\(M\models_A F\neg\mathsf{NZ}\); its counterexample query
is \(M\models_E G\mathsf{NZ}\)
\cite[\S2 and fig.~1]{BiereCimattiClarkeZhu1999}.

At the report's illustrative horizon \(k=2\), let
\(x_0,x_1,x_2\) be the three state-vector copies. The unfolding is
\[
 f_M=T(x_0,x_1)\land T(x_1,x_2).
\]
Its six bit constraints are, for \(i=0,1\),
\[
 x_{i+1}[0]=x_i[1],\qquad
 x_{i+1}[1]=x_i[2],\qquad x_{i+1}[2]=1.
\]
For \(j=0,1,2\), the possible closing edge is
\[
\begin{split}
 L_j=T(x_2,x_j)\iff{}&
 x_j[0]=x_2[1]\ \land\ x_j[1]=x_2[2]\\
 &{}\land\ x_j[2]=1.
\end{split}
\]
No-loop semantics cannot certify \(G\mathsf{NZ}\). Every
selected-loop translation at position zero requires nonzero
states throughout the displayed prefix, so the SAT formula is
\[
 f_M\land(L_0\lor L_1\lor L_2)
       \land\bigwedge_{j=0}^{2}\mathsf{NZ}(x_j).
\]
Setting all nine state bits to \(1\) satisfies every unfolding
constraint, every closing-edge alternative, and every nonzero
test. The resulting infinite execution is constantly \(111\),
and therefore refutes eventual zero
\cite[\S2, eq.~(1)]{BiereCimattiClarkeZhu1999}.

The horizon \(2\) illustrates unfolding; it is not minimal.
At \(k=0\), the same query is already
\[
 I(x_0)\land T(x_0,x_0)\land\mathsf{NZ}(x_0).
\]
The self-loop equations force all three bits to \(1\), which
satisfies this formula. Thus a search starting at zero finds
the counterexample immediately. The horizon counts unfolded
transitions, while the closing self-loop is checked separately.

\paragraph{Recurrence diameter and finite stopping bounds.}
Let \(S\ne\varnothing\) be finite and \(T\subseteq S\times S\).
A simple directed path is a finite sequence \(s_0,\ldots,s_m\) whose
consecutive states are connected by \(T\) and whose states are pairwise
distinct. Its length is its number \(m\) of edges. Define
\[
\begin{split}
 \operatorname{recdiam}(S,T)
 :=\max\{m\in\mathbb N:\;&
   \exists s_0,\ldots,s_m\in S,\\
 &\bigwedge_{i=0}^{m-1}T(s_i,s_{i+1})
   \ \land\
   \bigwedge_{0\leq i<j\leq m}s_i\ne s_j\}.
\end{split}
\]
A one-state path shows that the set is nonempty. A simple path has at most
\(|S|\) vertices, so
\[
 0\leq\operatorname{recdiam}(S,T)\leq |S|-1.
\]
Equivalently, \(d=\operatorname{recdiam}(S,T)\) is the least nonnegative
integer for which
\[
 \forall s_0,\ldots,s_{d+1}\in S,\qquad
 \bigwedge_{i=0}^{d}T(s_i,s_{i+1})
 \ \Longrightarrow\
 \bigvee_{0\leq i<j\leq d+1}s_i=s_j.
\]
To prove equivalence, a path of \(d+1\) edges cannot be simple by maximality
of \(d\). Conversely, if \(c<d\), the first \(c+2\) vertices of a longest
simple path violate the displayed property with \(c\) in place of \(d\).
The repeated pair may occur anywhere along the path.

This also gives a symbolic upper-bound certificate. For any \(r\geq0\),
unsatisfiability of
\[
 \left(\bigwedge_{i=0}^{r}T(s_i,s_{i+1})\right)
 \land
 \left(\bigwedge_{0\leq i<j\leq r+1}s_i\ne s_j\right)
\]
is equivalent to \(\operatorname{recdiam}(S,T)\leq r\), with all state
vectors constrained to \(S\). No initial-state condition is present:
this certificate bounds simple paths throughout the graph.

\paragraph{Reachability and invariant verification.}
Return to the finite total structure \(M=(S,I,T,L)\), and let \(p\) be a
Boolean state predicate. Put \(d=\operatorname{recdiam}(S,T)\). Then
\[
 M\models_E Fp
 \quad\Longleftrightarrow\quad
 \exists k\in\{0,\ldots,d\}\;(M\models_E^k Fp).
\]
If an initial execution reaches \(p\), choose a shortest initial-to-\(p\)
path. A repeated vertex would permit removal of the intervening cycle,
contradicting shortestness. Its length is therefore at most \(d\).
Totality extends it indefinitely, while its endpoint supplies the bounded
eventuality witness. The reverse implication follows from bounded soundness.

For every \(r\in\mathbb N\), define
\[
 C_r:=I(s_0)\land
       \bigwedge_{i=0}^{r-1}T(s_i,s_{i+1})
       \land\bigvee_{i=0}^{r}p(s_i).
\]
If \(r\geq d\), then \(M\models_A G\neg p\) exactly when \(C_r\) is
unsatisfiable.
Any satisfying assignment is a bad run prefix. Conversely, every reachable
bad state has a path of length at most \(d\leq r\); totality extends that
path to length \(r\), preserving the witnessed occurrence of \(p\).
Thus \(r=|S|-1\) is always a sufficient safety bound, and a smaller
recurrence-diameter certificate may improve it. If \(I=\varnothing\), every
\(C_r\) is unsatisfiable and the universal property holds vacuously.

Reachability may admit an even smaller bound. For \(I\ne\varnothing\), let
\(\operatorname{dist}(I,s)\) be the shortest initial-to-\(s\) distance for a
reachable state, and put
\[
 \delta_I=\max\{\operatorname{dist}(I,s):s\text{ is reachable from }I\}.
\]
The same shortest-path proof gives
\(\delta_I\leq d\), and \(C_{\delta_I}\) suffices for safety. The quantity
\(\delta_I\) measures shortest reachability distances; recurrence diameter
measures longest simple paths.

\paragraph{Certifying the ordinary diameter.}
Let \(\operatorname{dist}(u,v)\) be the shortest \(T\)-path
length from \(u\) to \(v\), defined for reachable pairs and
allowing length zero. The all-state diameter is
\[
 \Delta=\max\{\operatorname{dist}(u,v):
                   u,v\in S,\ v\text{ is reachable from }u\}.
\]
The set includes every \((u,u)\), so it is nonempty.
Shortest paths are simple; hence
\(0\leq\Delta\leq d\leq|S|-1\), and, when
\(I\ne\varnothing\), also \(\delta_I\leq\Delta\).

For any integer \(r\geq0\), the following quantified formula
is true exactly when \(\Delta\leq r\):
\[
\begin{split}
 \forall s_0,\ldots,s_{r+1}\in S\
 \exists t_0,\ldots,t_r\in S:\qquad&
 \bigwedge_{i=0}^{r}T(s_i,s_{i+1})
 \ \Rightarrow\\
 &\left(t_0=s_0\ \land\
     \bigwedge_{i=0}^{r-1}T(t_i,t_{i+1})
       \ \land\ \bigvee_{i=0}^{r}t_i=s_{r+1}\right).
\end{split}
\]
If \(\Delta\leq r\), the endpoint of any displayed
\((r+1)\)-edge path can be reached from its start within
\(r\) edges. Totality extends that shorter path to exactly
\(r\) edges if necessary, preserving its endpoint occurrence.
This supplies the existential vectors. Conversely, if a
shortest path had more than \(r\) edges, the formula would
shorten its first \(r+1\) edges to at most \(r\); following
the remaining suffix would contradict shortestness.
This proves both directions, including \(r=0\).
Under Boolean state encoding the formula is a quantified
Boolean formula with a universal-existential prefix,
distinct from the unsatisfiable-simple-path SAT certificate
for recurrence diameter
\cite[def.~19 and thms.~20--21]{BiereCimattiClarkeZhu1999}.

The ordinary diameter bounds unrestricted target reachability,
but not a path required to remain in an arbitrary predicate.
For example, take a chain \(a_0\to a_1\to a_2\to a_3\),
add a hub \(z\) with edges in both directions between it
and every \(a_i\), and give \(a_3\) a self-loop. The
unrestricted diameter is \(2\). If the until left operand
holds exactly at \(a_0,a_1,a_2\) and its right operand
exactly at \(a_3\), the shortest permitted witness from
\(a_0\) has three edges: the shortcut through \(z\) is
not allowed. The longest-simple-path bound \(d\) still
applies to this constrained witness.

\paragraph{A bounded witness for existential persistence.}
For the same state predicate \(p\),
\[
 M\models_E Gp
 \quad\Longleftrightarrow\quad
 \exists k\in\{0,\ldots,d\}\;(M\models_E^k Gp).
\]
An infinite initial \(p\)-path contains a cycle in the subgraph of
\(p\)-states. Choose a simple such cycle and a shortest initial path to
its vertex set within that subgraph. The stem is simple and has no cycle
vertex before its endpoint. Following the stem and then the cycle, but
omitting the repeated closing vertex, is consequently a simple path of
length \(k\leq d\). The closing edge gives a \(k\)-loop consisting entirely
of \(p\)-states. It is the required bounded witness. Conversely, bounded
soundness gives an infinite \(p\)-path. This argument concerns a state
predicate \(p\); an arbitrary temporal formula requires additional state
that records its acceptance obligations.

\paragraph{Padding a lasso to one fixed horizon.}
Suppose an execution has a stem of \(u\geq0\) states followed by
a repeating block of \(v\geq1\) states. Its first lasso encoding
has horizon \(k=u+v-1\), with loop entry \(u\).
For any \(K\geq k\), take its prefix through position \(K\) and put
\[
 \ell_K=K+1-v.
\]
Then \(u\leq\ell_K\leq K\), and periodicity gives
\(s_{K+1}=s_{\ell_K}\). Thus the actual next transition closes
the prefix back to \(\ell_K\). The loop
\(s_{\ell_K}\cdots s_K\) is a rotation of the same period,
so its infinite repetition after the retained prefix represents
the very same execution, not just a path with the same visited
states. On a selected loop, bounded semantics agrees with
ordinary LTL semantics. Consequently a satisfying lasso at
horizon \(k\) yields a satisfying SAT encoding at every
horizon \(K\geq k\), however deeply temporal operators are nested.

\paragraph{A completeness bound for arbitrary LTL.}
Let \(h\) be an LTL formula over a finite proposition set \(AP\), and let
\[
 \mathcal A_h=(Q,Q_0,\delta,F)
\]
be a finite nondeterministic B\"uchi automaton, without epsilon transitions,
whose accepted words are exactly the infinite label sequences satisfying
\(h\). Here \(Q\ne\varnothing\), \(Q_0\subseteq Q\),
\(\delta:Q\times2^{AP}\to\mathcal P(Q)\), and \(F\subseteq Q\);
acceptance requires infinitely many visits to \(F\).
Existence of this automaton is Theorem~\ref{thm:p09-ltl-buchi}; its
construction is given in Section~\ref{sec:p09-presentation-n-c-013}.

The product has vertices \(S\times Q\), initial vertices \(I\times Q_0\),
accepting vertices \(S\times F\), and edges
\[
 (s,q)\longrightarrow(t,q')
 \quad\Longleftrightarrow\quad
 T(s,t)\ \land\ q'\in\delta(q,L(s)).
\]
Thus \(q\) is the automaton state before reading the label of \(s\).
Let \(d_P\) be the recurrence diameter of this product graph. Then
\[
 M\models_E h
 \quad\Longleftrightarrow\quad
 \exists k\leq d_P\;(M\models_E^k h),
 \qquad d_P\leq |S||Q|-1.
\]
An accepting product run visits some accepting vertex infinitely often.
That vertex lies on a positive cycle; a shortest cycle through it is
simple apart from its repeated endpoint. Choose a shortest initial path
to that cycle's vertex set. The stem is simple and disjoint from the
cycle before entry. The stem followed by one cycle traversal, omitting
the closing vertex, is therefore a simple product path of length at most
\(d_P\). Repeating the cycle is accepting. Its projection onto \(S\)
is a \(k\)-loop satisfying \(h\), proving the forward implication.
The reverse implication follows from bounded soundness.

This proof uses one B\"uchi acceptance set \(F\). If an LTL translation
first produces a generalized B\"uchi automaton with several acceptance
sets, convert it to ordinary B\"uchi acceptance and include the added
acceptance-tracking states in \(Q\). The count \(|Q|\) is the size of that
actual automaton. Consequently, take \(h\) to be the negation-normal form of
\(\neg\varphi\) and construct \(\mathcal A_h\). The proof produces a
satisfying lasso of horizon at most \(d_P\) whenever the property fails.
By lasso padding, the universal property \(\varphi\) is therefore proved
by unsatisfiability of the single encoding for \(h\) at horizon
\(|S||Q|-1\), or at a certified smaller upper bound on \(d_P\).
The bound is finite but need not be small
\cite[\S5]{BiereCimattiClarkeZhu1999}.

\paragraph{Systems with uniformly bounded stems and periods.}
Suppose there are integers \(U\geq0\) and \(V\geq1\) such that
every initial infinite execution of \(M\) has a lasso
representation with at most \(U\) stem states and at most
\(V\) period states. For every LTL formula \(h\), put
\[
 K=U+V-1.
\]
Then
\[
 M\models_E h
 \quad\Longleftrightarrow\quad
 \llbracket M,h\rrbracket_K\text{ is satisfiable}.
\]
Indeed, any satisfying execution has the assumed representation
of horizon at most \(K\). Padding gives a satisfying encoding
at \(K\); the converse is bounded soundness. The bound depends
on the system's permitted executions and not on the size or
temporal depth of \(h\). This makes explicit the bounded-loop
case discussed in Section~5.2 \cite[\S5.2]{BiereCimattiClarkeZhu1999}.

For a deterministic total transition relation, let \(N=|S|\).
Following the unique successor from an initial state gives a
first repetition within \(N+1\) visited states. If it repeats
position \(i\) at position \(j\), \(0\leq i<j\leq N\),
determinism forces all subsequent states to repeat with
period \(j-i\). The stem and one period together therefore
contain \(j\leq N\) states. Regardless of which initial
state is selected, lasso padding gives the uniform bound
\[
 K=N-1
\]
for every LTL formula. This uses a bound on the sum of the
stem and period lengths; it does not add their separate
worst-case bounds. At \(N=1\), the self-loop at horizon
zero still supplies the successor positions needed to
interpret arbitrarily many nested next operators.

Finiteness alone does not give this property for every
nondeterministic execution. A graph with states \(0,1\)
and all four transitions can follow
\(0\,1\,00\,1\,000\,1\,0000\,1\cdots\), which is not
ultimately periodic: its visits to \(1\) have unbounded
gaps. General finite-state LTL verification still has
lasso witnesses through the automaton-product argument,
but their uniform bound may depend on the formula's
acceptance obligations.

\paragraph{A linear witness bound for the \(F,G\) fragment.}
Suppose \(h\) is built from constants, literals, conjunction,
disjunction, and only the temporal operators \(F,G\).
Let \(m\) count its distinct temporal subformulas and put
\(N=|S|\geq1\). A sufficient bound is
\[
 K=2(m+1)N-2,\qquad
 M\models_E h
 \quad\Longleftrightarrow\quad
 \llbracket M,h\rrbracket_K\text{ is satisfiable}.
\]
The following construction proves an explicit linear
bound of the kind considered in Section~5.2; it does
not assert that this coefficient is optimal
\cite[\S5.2]{BiereCimattiClarkeZhu1999}.

\emph{Divide a witness into intervals.}
Assume an initial infinite path \(\pi\) satisfies \(h\).
Along that path, the truth of any \(F\theta\) can change
only from true to false, since later suffixes have fewer
possible future witnesses. The truth of \(G\theta\)
can change only from false to true. Record the truth
values of all \(m\) temporal subformulas at each position.
There are at most \(m+1\) consecutive intervals with a
constant record, the last of which is infinite.
Within each finite interval, erase cycles from its
state sequence while keeping the first and last states.
If these coincide, one occurrence can serve both roles.
At most \(N\) states remain, and each retains the interval's
record. The original edge across every interval boundary
still connects the retained endpoint states.

\emph{Build a short recurring suffix.}
Let \(Z\) be the nonempty set of states occurring infinitely
often in the final interval. Beyond some point, only
states of \(Z\) occur. Between sufficiently late occurrences
of any two of its states the original path stays in \(Z\),
so its induced graph is strongly connected. Every state
in \(Z\) also has a positive closed walk.

For every subformula \(F\theta\) whose final recorded
value is true, \(\theta\) must hold infinitely often:
otherwise \(F\theta\) would eventually be false. Choose
a recurring state witnessing \(\theta\).
For every \(G\theta\) whose final value is false, choose
a recurring state witnessing \(\neg\theta\), which must
likewise occur infinitely often. There are at most \(m\)
such obligations. Once a temporal record is fixed, the
value of a subformula is determined by that record
and the state's atomic labels, so these choices give
the required values at all sufficiently late occurrences
of the chosen states.

Connect the selected witnesses cyclically by shortest
paths inside \(Z\). With \(r\geq1\) witnesses the total
length is at most \(r(N-1)\), unless all connecting
paths have length zero. In that case use a shortest
positive cycle through their common state, of length
at most \(N\). With no obligations, choose any simple
positive cycle in \(Z\). Thus there is a nonempty
closed walk of length
\[
 b\leq\max\{1,m\}N
\]
that supplies every recurring obligation.
From the start of the final interval, follow the
original path to its chosen cycle-entry state and
erase cycles. This contributes at most \(N\) vertices.
Together with the retained finite intervals, the stem,
excluding its cycle-entry vertex, has at most
\[
 a\leq(m+1)N-1
\]
states. Append infinite repetition of the chosen closed
walk. Every edge is an edge of \(M\).

\emph{Preserve temporal truth.}
Assign each retained stem position its interval record
and every new cycle position the final record. Structural
induction shows that the recorded values are its actual
subformula values on the new path. Atomic and Boolean
clauses follow from the retained state labels. For
\(F\theta\) recorded false, every later original position
satisfies \(\neg\theta\). This covers both the retained
later stem positions and the recurring states used by
the new loop; induction therefore makes the eventuality
false on the new path.

For \(F\theta\) recorded true, if its final value is true,
the loop supplies a \(\theta\)-witness. Otherwise,
\(\theta\) holds at the last position before that
eventuality changes to false. That position is an
interval endpoint, whose state and record were kept.
Induction supplies the corresponding future witness.
For \(G\theta\) recorded true, every later original
position satisfies \(\theta\), and induction preserves
that fact at all retained stem and loop positions.
For \(G\theta\) recorded false, either its value stays
false and the loop supplies a \(\neg\theta\)-witness,
or it changes to true. In the latter case \(\theta\)
is false at the last position before the change,
again a retained interval endpoint.
These four cases prove the induction, including nested
\(F,G\) operators. In particular the new initial path
still satisfies \(h\).

Its natural lasso horizon satisfies
\[
 a+b-1
 \leq (m+1+\max\{1,m\})N-2
 \leq 2(m+1)N-2.
\]
Padding gives a witness at exactly \(K\); bounded
soundness proves the converse implication of the theorem.
At \(m=0,N=1\), this gives \(K=0\), with the required
self-loop. Next is excluded from the fragment because
the deletion of intermediate states does not in general
preserve its meaning.

\paragraph{Branching certificates for existential CTL.}
\label{sec:p09-n-c-017-biere-ectl-certificates}
For existential CTL, different subformulas may require
different paths. The state-certificate clauses below are a
handbook construction, separate from the report's unprinted
bounded ECTL semantics. They use the state-formula grammar
\[
 \theta::=\top\mid\bot\mid p\mid\neg p
 \mid\theta\land\theta\mid\theta\lor\theta
 \mid EX\theta\mid E[\theta\,U\,\theta]\mid EG\theta.
\]
Negation occurs only at atoms. Ordinary state semantics
is the CTL semantics already given: until has a finite
target-reaching prefix with the left operand before
the target, and \(EG\) has an infinite witness through
states satisfying its operand. Define
\(EF\alpha:=E[\top\,U\,\alpha]\).
Existential release is also expressible:
\[
 E[\alpha\,R\,\beta]
 \equiv EG\beta\lor E[\beta\,U\,(\alpha\land\beta)].
\]
The two alternatives are respectively an infinite
\(\beta\)-path or an \(\alpha\)-witness with \(\beta\)
holding through its endpoint.

Write \(M,s\Vdash_k\theta\) for the following certificate
relation, with \(k\geq0\). Its constants and literals
have their ordinary state meanings; its Boolean clauses
use \(\Vdash_k\) on the children. The temporal clauses are
\[
\begin{aligned}
 M,s\Vdash_k EX\alpha
 \quad\Longleftrightarrow\quad&
 k\geq1\ \land\
 \exists t\in S\,[T(s,t)\land M,t\Vdash_k\alpha],\\
 M,s\Vdash_k E[\alpha\,U\,\beta]
 \quad\Longleftrightarrow\quad&
 \exists r\in\{0,\ldots,k\}\ \exists s_0,\ldots,s_r\in S:\\
 &s_0=s\ \land\
   \bigwedge_{j=0}^{r-1}T(s_j,s_{j+1})\\
 &{}\land M,s_r\Vdash_k\beta\
   \land\bigwedge_{j=0}^{r-1}M,s_j\Vdash_k\alpha,\\
 M,s\Vdash_k EG\alpha
 \quad\Longleftrightarrow\quad&
 \exists r\in\{0,\ldots,k\}\
 \exists\ell\in\{0,\ldots,r\}\
 \exists s_0,\ldots,s_r\in S:\\
 &s_0=s\ \land\
   \bigwedge_{j=0}^{r-1}T(s_j,s_{j+1})\
   \land T(s_r,s_\ell)\\
 &{}\land\bigwedge_{j=0}^{r}M,s_j\Vdash_k\alpha.
\end{aligned}
\]
The horizon is local to each temporal certificate;
recursive calls retain \(k\) but use proper subformulas.
The definition is consequently well founded by formula
size, not by decrementing the horizon. Each child
occurrence can have its own witnesses.
At \(k=0\), \(EX\) is false, until requires its right
operand at the starting state, and \(EG\) requires a
self-loop and its operand certificate.
This explicitly defined state-certificate relation
is distinct from the earlier bounded LTL path relation.

For every \(k\), certificates are sound:
\[
 M,s\Vdash_k\theta\quad\Longrightarrow\quad M,s\models\theta.
\]
Prove this simultaneously at all states by structural
induction. Next supplies an actual edge to a true child.
Until supplies its finite true prefix and target;
totality extends it to an infinite path. The \(EG\)
certificate supplies a lasso whose every state satisfies
the operand by induction. Positive Boolean combinations
preserve the implication.

Let \(d=\operatorname{recdiam}(S,T)\). For every
\(k\geq\max\{1,d\}\), the converse holds as well.
The next case has a one-edge witness. For until,
choose a shortest witness satisfying its state constraints.
A repeated state would permit cycle deletion without
losing those constraints, so its length is at most \(d\).
For \(EG\), choose a simple cycle within the operand's
state set and a shortest path from the given state
to that cycle's vertex set inside the same set.
The stem and cycle are disjoint before entry; omitting
the repeated closing vertex leaves a simple path of
length at most \(d\). In each case the induction
hypothesis supplies all child certificates at the
same \(k\). Thus
\[
 M,s\models\theta
 \quad\Longleftrightarrow\quad M,s\Vdash_k\theta
 \qquad(k\geq\max\{1,d\}).
\]
Increasing \(k\) only enlarges the permitted witness
families, so it also preserves certificate truth.
The positive lower bound \(1\) accounts for the explicit
next clause when \(d=0\).

For Boolean state encodings, compile these clauses
recursively to propositional formulas. Next introduces
a fresh successor vector; until disjoins over the
possible lengths and introduces its path vectors;
\(EG\) additionally disjoins over loop entries and
checks the closing edge. Each child occurrence gets
private fresh witness variables. All witness quantifiers
are existential and positive, so after renaming they
can be moved to the front and left implicit in a
SAT query. Structural induction gives exact correspondence
with \(\Vdash_k\). This is a correctness construction;
naive expansion is not claimed to have polynomial size.

Conjoin \(I(s)\) and existentially choose the starting
state to search for an initial ECTL witness. Negating
an ACTL property into this existential negation-normal
form gives its initial counterexample query, which is
complete at the displayed bound. The construction
makes the branching required by existential CTL
explicit \cite[\S5.1]{BiereCimattiClarkeZhu1999}.
For example, \(EXp\land EXq\) can hold at a state
with one \(p\)-only successor and one \(q\)-only successor,
although no single successor satisfies \(p\land q\).
Requiring the two conjuncts to share a path would
change the property.

\paragraph{The implementation and positive verification experiments.}
\label{sec:p09-n-c-017-biere-positive-experiments}
The prototype \textsc{BMC} accepts a subset of SMV's input
language and can emit an SMV program or propositional formulas.
Its SAT outputs support DIMACS input for \textsc{SATO}, based on
the Davis--Putnam procedure, and the input format of
\textsc{PROVE}, based on St\aa lmarck's method
\cite[\S6]{BiereCimattiClarkeZhu1999}.

One benchmark checks a sequential shift-and-add multiplier
against the combinational C6288 reference on the same input
words. The operands have sixteen bits each; the experiment
retains sixteen result bits and an overflow bit, checking the
corresponding outputs individually when the sequential
multiplier finishes. A second benchmark rotates a register
file and verifies an invariant relating it to another file:
equal contents in two registers imply equal contents in their
respective neighbors. The authors also checked a finite bound
for that rotation model. These experiments show bounded SAT
reasoning used for complete verification when its stopping
bound is established, as well as for finding counterexamples.
Their timings compare particular encodings, solver settings,
and BDD implementations; they are not a worst-case complexity
improvement \cite[\S6 and tables~1--2]{BiereCimattiClarkeZhu1999}.

\paragraph{Interpreting a bounded experiment.}
In the distributed-mutual-exclusion experiment,
the liveness claim that every request is eventually acknowledged depends on a
fairness constraint for each asynchronous gate.  Because no maximum
counterexample length was known, the authors could not verify the correct
design completely; they established only the absence of counterexamples at
\(k=5\) and \(k=10\).  Removing fairness constraints from the buggy design
produced a two-loop counterexample immediately.  This source therefore
supports fast bounded bug finding and a completeness theorem quantified over
\(k\), not an unbounded proof from one finite \(\mathsf{UNSAT}\) result
\cite[\S\S5--6 and tables~3--4]{BiereCimattiClarkeZhu1999}.

\subsubsection{Clarke et al.'s counterexample-guided refinement}
\label{sec:p09-n-c-017-cegar}

The source is the 43-page \emph{Journal of the ACM} article, pp.~752--794, not
merely its preliminary CAV~2000 version.
Its first page carries ACM's conditional personal/classroom copying notice,
not an open-content licence.  The reconstruction below states the mathematical
objects and proof boundaries without reproducing the article
\cite{ClarkeEtAl2003}.

\paragraph{Concrete structures, ACTL*, and simulation.}
In the notation used for the construction, a concrete Kripke structure is
\[
 M=(S,I,R,L),
\]
where \(S\) is finite, \(\varnothing\ne I\subseteq S\),
\(R\subseteq S\times S\) is total, and \(L:S\to2^{\mathcal A}\).  A path is
an infinite sequence \(s_0s_1\ldots\) with \(R(s_i,s_{i+1})\) at every
position.  ACTL* is the fragment of CTL* in which path quantification is
universal and negation occurs only before atomic formulas.  The restriction is
essential: the preservation result below is not a theorem for arbitrary CTL*.
The component order \(M=(S,I,R,L)\) is used throughout this treatment.

For Kripke structures \(M\) over \(\mathcal A\) and \(M'\) over
\(\mathcal A'\subseteq\mathcal A\), a relation
\(H\subseteq S\times S'\) is a simulation relation when, for every
\((s,s')\in H\),
\[
\begin{aligned}
 L(s)\cap\mathcal A'&=L'(s'),\\
 R(s,t)&\Longrightarrow
   \exists t'\in S'\;\bigl(R'(s',t')\land(t,t')\in H\bigr).
\end{aligned}
\]
The structure \(M'\) simulates \(M\), written \(M\preceq M'\), when such an
\(H\) also relates every concrete initial state to some abstract initial state.
Theorem~2.3 states the one-way preservation law
\[
 M\preceq M'\ \land\ M'\models\varphi
 \quad\Longrightarrow\quad
 M\models\varphi
 \qquad(\varphi\in\mathrm{ACTL}^*).
\]

The proof is simultaneous induction on ACTL* state and path formulas.
The atomic case is the displayed label equality; conjunction and disjunction
use the induction hypothesis.  From a related pair of states, the transition
clause lifts each concrete path step by step to a pointwise related abstract
path.  If every abstract path satisfies a path formula, the lifted path does;
the clauses for next, until, release, eventually, and always follow from the
induction hypotheses at corresponding positions.  Atomic-only negation is
preserved by label equality.  Universal quantification then yields the state
clause.  Existential path quantification fails for a different reason: an
abstract existential witness can use added behaviour and need not have a
concrete realization.  Unrestricted negation can reverse a preservation
implication by taking complements.  ACTL* excludes both failure mechanisms
\cite[\S2.4 and thm.~2.3]{ClarkeEtAl2003}.

\paragraph{Existential abstraction and appropriateness.}
An abstraction function is a surjection
\[
 h:S\twoheadrightarrow\widehat S.
\]
It induces \(d\equiv_h e\) exactly when \(h(d)=h(e)\).  The abstract structure
\(\widehat M_h=(\widehat S,\widehat I_h,\widehat R_h,\widehat L_h)\) is
defined componentwise by
\[
\begin{aligned}
 \widehat I_h(\widehat d)
  &\iff\exists d\in S\;(h(d)=\widehat d\land d\in I),\\
 \widehat R_h(\widehat d_1,\widehat d_2)
  &\iff\exists d_1,d_2\in S\;
       \bigl(h(d_1)=\widehat d_1\land h(d_2)=\widehat d_2
             \land R(d_1,d_2)\bigr),\\
 \widehat L_h(\widehat d)
  &=\bigcup_{h(d)=\widehat d}L(d).
\end{aligned}
\]
The existential transition clause is the exact source of added behaviour: two
successive abstract edges may be witnessed by incompatible concrete states in
their shared block.

For a specification \(\varphi\), the paper calls \(h\) appropriate when every
atomic subformula of \(\varphi\) is constant on each \(\equiv_h\)-class:
\[
 d\equiv_h e
 \quad\Longrightarrow\quad
 \bigl(d\models f\iff e\models f\bigr)
 \quad\text{for every atomic }f\text{ in }\varphi.
\]
On the proposition set occurring in \(\varphi\), appropriateness makes the
union label of an abstract block equal to the common label of each represented
concrete state.  The graph
\[
 H_h=\{(s,h(s)):s\in S\}
\]
is then a simulation relation, because every concrete transition itself
witnesses the corresponding abstract transition and concrete initial states
map into \(\widehat I_h\).  Lemma~3.2 and Theorems~2.3 and~3.3 therefore give
\[
 \widehat M_h\models\varphi
 \quad\Longrightarrow\quad
 M\models\varphi.
\]
The converse is not valid: existential abstraction can add paths
\cite[\S3 and thm.~3.3]{ClarkeEtAl2003}.

\paragraph{Exact abstraction versus symbolic approximations.}
Computing the exact existential transition relation \(\widehat R_h\) can itself
be prohibitively expensive.  The article therefore defines an
\emph{approximation} to be the Kripke structure
\[
 \widetilde M=(\widehat S,\widetilde I,\widetilde R,\widehat L_h)
 \quad\text{such that}\quad
 \widehat I_h\subseteq\widetilde I
 \ \text{ and }\
 \widehat R_h\subseteq\widetilde R.
\]
This term names an overapproximating structure on the same abstract state set,
not another abstraction function.  Its necessity is computational, and it is
nonredundant with \(\widehat M_h\): the latter contains exactly the existentially
induced initial states and transitions, whereas \(\widetilde M\) may add more.
The identity relation simulates \(\widehat M_h\) in \(\widetilde M\), so
Theorem~3.5 gives
\[
 M\preceq\widehat M_h\preceq\widetilde M.
\]
Consequently, a successful ACTL* check on \(\widetilde M\) is still sound.
A transition in
\(\widetilde R\setminus\widehat R_h\) is a \emph{virtual transition}: unlike a
spurious path already possible in exact existential abstraction, it was added
by this further approximation.  Both kinds must later be checked against the
concrete relation \(R\).

The source applies existential abstraction to individual Boolean blocks before
conjoining them.  The block notation must be typed before that rule can be
used.  For an initial-predicate block \(B\) and a transition-relation block
\(Q\), define their existential images by
\[
\begin{aligned}
 (B)_h(\widehat s)
   &\iff \exists s\in S\;
      \bigl(h(s)=\widehat s\land B(s)\bigr),\\
 (Q)_h(\widehat s,\widehat t)
   &\iff \exists s,t\in S\;
      \bigl(h(s)=\widehat s\land h(t)=\widehat t
             \land (s,t)\in Q\bigr).
\end{aligned}
\]
Thus \((B)_h\) is a predicate on \(\widehat S\), whereas \((Q)_h\) is a
relation on \(\widehat S\times\widehat S\); neither is a new abstraction
function.  If the Boolean encodings are partitioned as
\[
 I=I_1\land\cdots\land I_r,
 \qquad
 R=R_1\land\cdots\land R_q,
\]
the source's \emph{early approximation} is exactly
\[
\begin{aligned}
 \mathcal T(I)
   &=\{\widehat s\in\widehat S:
       \bigwedge_{\ell=1}^{r}(I_\ell)_h(\widehat s)\},\\
 \mathcal T(R)
   &=\{(\widehat s,\widehat t)\in\widehat S^2:
       \bigwedge_{j=1}^{q}(R_j)_h(\widehat s,\widehat t)\}.
\end{aligned}
\]
These definitions make the initial component as well as the transition
component eliminable.  The full structure is
\[
 \widetilde M_h^{early}
   =(\widehat S,\mathcal T(I),\mathcal T(R),\widehat L_h).
\]
Because
\[
 \widehat I_h\subseteq\mathcal T(I),
 \qquad
 \widehat R_h\subseteq\mathcal T(R),
\]
this is a full approximation in the preceding sense and
\[
 M\preceq\widehat M_h\preceq\widetilde M_h^{early}.
\]
Abstracting the blocks before conjoining them can be cheaper, but may add virtual
transitions---and, if the first inclusion is strict, extra abstract initial
states---that exact abstraction of the conjunction would exclude
\cite[\S3.1 and thm.~3.5; pp.~763--765]{ClarkeEtAl2003}.

The source's smallest example is a traffic light with the concrete cycle
\[
 red\longrightarrow green\longrightarrow yellow\longrightarrow red.
\]
Let \(h(red)=\widehat{red}\) and
\(h(green)=h(yellow)=\widehat{go}\).  The concrete system satisfies
\(\mathsf{AG\,AF}(state=red)\), but the abstract block
\(\widehat{go}\) has a self-loop witnessed by \(green\to yellow\).  Hence the
abstract trace
\(\widehat{red}\,\widehat{go}^{\,\omega}\) violates the formula.  Because no
concrete execution can remain forever among \(green,yellow\), that trace is
spurious.  This example exhibits the failure of the converse preservation
direction rather than a failure of Theorem~3.3
\cite[ex.~3.4 and fig.~2]{ClarkeEtAl2003}.

\paragraph{Guarded programs and their transition relation.}
A guarded finite-state program \(P\) has a finite state-variable set and
nonempty finite variable domains:
\[
 V=\{v_1,\ldots,v_n\},
 \qquad
 \varnothing\ne D_v\quad(v\in V),
 \qquad
 S=\prod_{v\in V}D_v.
\]
Each variable \(v_i\) has a nonempty initial-value set \(I_i\subseteq D_{v_i}\)
and an ordered list of guarded updates
\[
 (C_i^1,A_i^1),\ldots,(C_i^{k_i},A_i^{k_i}),\qquad k_i\geq1.
\]
Each \(C_i^j:S\to\{\mathsf{false},\mathsf{true}\}\) is a predicate on the
current state. Each \(A_i^j:S\to
\mathcal P(D_{v_i})\setminus\{\varnothing\}\) gives the permitted next values.
A deterministic expression denotes its singleton value set; a finite
set-valued assignment permits any member. All expressions are interpreted
by total functions of the declared finite types.

Require at least one guard in each list to hold at every state. An explicit
last guard \(\mathsf{true}\) supplies a default when necessary.
The effective guard for branch \(j\) is
\[
 G_i^j(s):=C_i^j(s)\land
                \bigwedge_{\ell=1}^{j-1}\neg C_i^\ell(s).
\]
Exactly one effective guard holds. Thus the first applicable branch,
not an arbitrary true branch, supplies the update. Define
\[
\begin{aligned}
 T_i(s,t)&\iff
   \bigvee_{j=1}^{k_i}
       \bigl(G_i^j(s)\land t(v_i)\in A_i^j(s)\bigr),\\
 I&=\{s\in S:\forall i\in\{1,\ldots,n\}\;s(v_i)\in I_i\},\\
 R(s,t)&\iff\bigwedge_{i=1}^{n}T_i(s,t).
\end{aligned}
\]
All right-hand sides are evaluated in the same current state \(s\).
The conjunction implements simultaneous updates; one block does not read
another block's newly chosen next value. For example,
\(A_i^1(s)=\{s(v_j)\}\) with guard \(\mathsf{true}\) copies the current value
of \(v_j\) when \(D_{v_j}\subseteq D_{v_i}\), while
\(A_i^1(s)=\{0,1\}\) for a Boolean variable permits both
next values independently.

The initial set is nonempty because every \(I_i\) is nonempty. For each
current state, choose one permitted value from the unique effective branch
of each variable. The resulting tuple is a successor, proving that \(R\)
is total. At \(n=0\), the state space is the singleton empty tuple, both
empty conjunctions are true, and the transition is its self-loop. This gives
the required finite total Kripke structure for the program
\cite[\S\S2.2, 4.2.1]{ClarkeEtAl2003}.

\paragraph{The initial abstraction.}
Let \(\operatorname{Atoms}(P)\) contain the atomic formulas in the
specification and in the control-flow conditions of the program's transition
blocks, and label each state by
\[
 L(s)=\{a\in\operatorname{Atoms}(P):s\models a\}.
\]
Here satisfaction of a program predicate evaluates its expressions on the
current variable assignment and combines the atomic truth values by the
usual Boolean clauses. Definition~4.1 says that, for any formula set \(F\),
\[
 d\equiv_F e
 \iff
 d,e\in S\ \land\
 \forall f\in F\;\bigl(d\models f\iff e\models f\bigr).
\]
The corresponding abstraction is the quotient map
\[
 h_F:S\twoheadrightarrow S/{\equiv_F},
 \qquad
 h_F(d)=[d]_{\equiv_F}.
\]
The map definition is necessary: the kernel alone determines a quotient only
up to renaming of its classes.  If
\(F\supseteq\operatorname{Atoms}(\varphi)\), Proposition~4.2 proves that
\(h_F\) is appropriate: a distinguishing specification atom would itself be
an element of \(F\), contradicting \(F\)-equivalence.  Definition~4.3 then
sets
\[
 init:=h_{\operatorname{Atoms}(P)}.
\]
The source chooses the program's control atoms as well as the specification
atoms so that the first abstract model retains the control-flow distinctions
visible in the program.  Proposition~4.2 also makes clear that this choice is a
design decision, not a necessary new semantic primitive
\cite[defs.~4.1, 4.3 and prop.~4.2; pp.~768--769]{ClarkeEtAl2003}.

\paragraph{The product representation required by symbolic refinement.}
The later refinement algorithm operates on coordinate equivalences, so it
requires more structure than the arbitrary surjection \(h_F\).  Suppose
\[
 S=D_1\times\cdots\times D_m,\qquad
 h_j:D_j\twoheadrightarrow\widehat D_j,
\]
and define
\[
\begin{aligned}
 h&=(h_1,\ldots,h_m),\\
 h(d_1,\ldots,d_m)
   &=(h_1(d_1),\ldots,h_m(d_m)),\\
 a\equiv_j b&\iff h_j(a)=h_j(b).
\end{aligned}
\]
Then \(h\) is a \emph{composite abstraction function}, a restricted kind of
abstraction function whose standing is a product map, not a new state or
transition relation.  Its full equivalence is eliminable:
\[
 (d_1,\ldots,d_m)\equiv_h(e_1,\ldots,e_m)
 \iff
 \bigwedge_{j=1}^{m}d_j\equiv_j e_j.
\]
The distinction is nonredundant because a general \(h_F\) can couple
coordinates and need not factor this way; the later componentwise refinement
operation would then be undefined.

The article constructs the required product representation of \(init\) rather
than assuming it.  For an atomic formula \(f\), let
\(\operatorname{var}(f)\) be the variables occurring in \(f\), and put
\[
 V_{\mathrm{obs}}
   =\bigcup_{f\in\operatorname{Atoms}(P)}\operatorname{var}(f).
\]
The source grammar permits an atom with
\(\operatorname{var}(f)=\varnothing\).  Its truth value is independent of the
state: every \(d,e\in S\) agree on all zero variables, so
\(d\models f\iff e\models f\).  Its conjunct in \(\equiv_{init}\) is therefore
tautological.  Such atoms remain in \(\operatorname{Atoms}(P)\) and in the
labelling, but are omitted when forming the product coordinates because they
cannot induce a variable coordinate and cannot distinguish two states.
This omission preserves \(\equiv_{init}\) exactly
\cite[\S4.2.1, p.~767; \S4.2.3, p.~770]{ClarkeEtAl2003}.

On the remaining atoms, two atoms \emph{interfere} exactly when their variable
sets intersect.  A \emph{formula cluster} is an equivalence class of the
reflexive-transitive closure of that interference relation on
\[
 \{f\in\operatorname{Atoms}(P):
     \operatorname{var}(f)\ne\varnothing\}.
\]
Two
variables in \(V_{\mathrm{obs}}\) are equivalent when they occur in atoms in
the same formula cluster; the resulting equivalence classes are the
\emph{variable clusters}.  Thus formula clusters partition precisely the
variable-bearing atoms, variable clusters partition \(V_{\mathrm{obs}}\), and
neither name denotes an additional semantic object.

For each unobserved variable \(v\in V\setminus V_{\mathrm{obs}}\), adjoin the
separate coordinate and empty formula family
\[
 VC_v=\{v\},
 \qquad
 FC_v=\varnothing.
\]
Separate coordinates are required here because the later refinement is
componentwise; grouping all unobserved variables into one coordinate would
permit coupled refinements not available when the variables are kept
independent. After adjoining these product factors and reindexing, write
\[
 FC_1,\ldots,FC_m,\qquad
 VC_1,\ldots,VC_m,\qquad
 D_j=\prod_{v\in VC_j}D_v.
\]
Let \(\pi_{VC_j}:S\to D_j\) be coordinate restriction to \(VC_j\).  Full-state
satisfaction is typed on \(S\), so a projected tuple needs an explicitly
restricted evaluation.  For \(f\in FC_j\), define the Boolean map
\(f^{(j)}:D_j\to\{\bot,\top\}\) by
\[
 f^{(j)}(x)=
 \begin{cases}
  \top,&d\models f,\\
  \bot,&d\not\models f,
 \end{cases}
 \qquad
 \text{for any }d\in S\text{ with }\pi_{VC_j}(d)=x.
\]
Such an extension \(d\) exists because every \(D_v\) is nonempty.  The value
is independent of the extension: if two full states have the same
\(VC_j\)-projection, they agree on every variable in
\(\operatorname{var}(f)\subseteq VC_j\), and hence give \(f\) the same truth
value.  Thus \(f^{(j)}\) is only the restriction of concrete formula
evaluation, introduced because the component quotient below uses it; it is not
a new satisfaction relation.

Define the component equivalence, quotient, and map by
\[
 x\equiv_j y
 \iff
 x,y\in D_j\ \land\
 \forall f\in FC_j\;
    \bigl(f^{(j)}(x)=f^{(j)}(y)\bigr),
 \qquad
 \widehat D_j=D_j/{\equiv_j},
 \qquad
 h_j:D_j\twoheadrightarrow\widehat D_j,
 \quad h_j(x)=[x]_{\equiv_j}.
\]
For each empty family the universal condition is vacuous, so
its equivalence is universal, its quotient is a singleton, and its component
map is constant.  The product equivalence therefore still ignores precisely
the coordinates ignored by \(init\), while retaining separate coordinates for
possible later refinements.  All components are actual surjections with
specified codomains, not merely kernels.
Theorem~4.8 states
\[
 \widehat M_h\cong\widehat M_{init}.
\]
Indeed,
\[
\begin{aligned}
 d\equiv_{init}e
 &\iff
   \forall f\in\operatorname{Atoms}(P)\;
      (d\models f\iff e\models f)\\
 &\iff
   \forall f\in\operatorname{Atoms}(P)\text{ with }
      \operatorname{var}(f)\ne\varnothing\;
      (d\models f\iff e\models f)\\
 &\iff
   \bigwedge_{j=1}^{m}
     \pi_{VC_j}(d)\equiv_j\pi_{VC_j}(e)\\
 &\iff d\equiv_he.
\end{aligned}
\]
Equal equivalence classes induce isomorphic exact abstract structures, which
is the promised bridge from the extensional \(init\) definition to the
component representation used below.

Example~4.7 makes the construction concrete.  For
\(x,y\in\{0,1,2\}\), Boolean \(reset\), and
\[
 \operatorname{Atoms}(P)
  =\{reset=TRUE,\ x=y,\ x<y,\ y=2\},
\]
the formula clusters are
\(\{x=y,x<y,y=2\}\) and \(\{reset=TRUE\}\), inducing variable clusters
\(\{x,y\}\) and \(\{reset\}\).  The \(x,y\)-component has the five classes
\[
\begin{gathered}
 \{(0,0),(1,1)\},\quad \{(0,1)\},\quad
 \{(0,2),(1,2)\},\\
 \{(1,0),(2,0),(2,1)\},\quad \{(2,2)\},
\end{gathered}
\]
while the \(reset\)-component is the identity
\cite[\S4.2.3, ex.~4.7, and thm.~4.8; pp.~769--771]{ClarkeEtAl2003}.

\paragraph{Combining exact and early transition abstraction.}
Let
\[
 R=T_1\cap\cdots\cap T_q,\qquad T_j\subseteq S\times S,\quad q\geq1.
\]
The blocks may correspond to variables or to groups of variables. They are
relations on full states, so a block may depend on current variables outside
the group whose next values it constrains. Fix a selection
\(J\subseteq\{1,\ldots,q\}\). The selected blocks share one concrete witness;
the remaining blocks are abstracted separately. Define
\[
\begin{split}
 \widetilde R_{h,J}(\widehat s,\widehat t)
 \iff{}&
 \exists s,t\in S\,
 \left(h(s)=\widehat s\land h(t)=\widehat t
       \land\bigwedge_{j\in J}T_j(s,t)\right)\\
 &{}\land
 \bigwedge_{j\notin J}(T_j)_h(\widehat s,\widehat t).
\end{split}
\]
Every term in this formula has the abstract domain
\(\widehat S\times\widehat S\); \((T_j)_h\) is the existential image already
defined above. Equivalently, the formula is
\[
 \left(\bigcap_{j\in J}T_j\right)_h
       \cap\bigcap_{j\notin J}(T_j)_h,
\]
where the empty concrete intersection is \(S\times S\).
Because \(h\) is surjective, its image is
\(\widehat S\times\widehat S\).
This is a precise combination of exact and early abstraction
\cite[\S\S3.1, 6]{ClarkeEtAl2003}.

The defining inclusions are
\[
 \widehat R_h\subseteq\widetilde R_{h,J}\subseteq\mathcal T(R).
\]
For the first inclusion, a concrete edge witnessing \(\widehat R_h\)
satisfies every \(T_j\), hence witnesses both the selected conjunction and
each unselected image. For the second, a witness satisfying all selected
blocks is separately a witness for each of their images. The other images
are explicit conjuncts. This proves both inclusions.
The extremes are
\[
 \widetilde R_{h,\varnothing}=\mathcal T(R),\qquad
 \widetilde R_{h,\{1,\ldots,q\}}=\widehat R_h.
\]
Moreover, if \(J\subseteq K\), then
\(\widetilde R_{h,K}\subseteq\widetilde R_{h,J}\): the shared \(K\)-witness
also witnesses \(J\), and every block in \(K\setminus J\) has its required
separate image. Enlarging the selection therefore never adds transitions.

Use the exact initial image and form the full structure
\[
 \widetilde M_{h,J}
   =(\widehat S,\widehat I_h,\widetilde R_{h,J},\widehat L_h).
\]
It has a nonempty initial set and a total transition relation, because
the exact abstract structure does and its transitions are included.
For an appropriate \(h\), the concrete-to-abstract simulation and the
identity simulations on \(\widehat S\) give
\[
 M\preceq\widehat M_h
   \preceq\widetilde M_{h,J}
   \preceq\widetilde M_h^{early}.
\]
Thus satisfaction of the ACTL* specification in
\(\widetilde M_{h,J}\) proves satisfaction in \(M\).
The selection \(J\) affects precision and computational cost, but every
choice has the same proved direction of preservation.

For example, let states be Boolean triples \((x,y,z)\), let
\(h(x,y,z)=(x,z)\), and put
\[
 T_1(s,t)\iff t(x)=s(y),\qquad
 T_2(s,t)\iff t(z)=s(y),\qquad
 T_3(s,t)\iff t(y)=s(y).
\]
With initial states \(\{(0,0,0),(0,1,0)\}\), every reachable state satisfies
\(x=z\). The exact abstract relation permits exactly the targets
\((x',z')\) with \(x'=z'\). Each separate image \((T_j)_h\), however, is
universal: its hidden \(y\)-witness can be chosen independently.
Early abstraction therefore permits unequal next values.
Selecting \(J=\{1,2\}\) forces the first two blocks to use the same current
\(y\), recovers \(x'=z'\), and gives
\(\widetilde R_{h,J}=\widehat R_h\). This exhibits precisely the information
retained by a shared witness.

\paragraph{The refinement loop.}
The exact method to which Theorem~4.24 belongs starts from the composite
representative \(h\) of \(init\) and repeats exactly three decisions.
\begin{enumerate}
\item Model-check \(\widehat M_h\) against \(\varphi\).  If it satisfies the
formula, Theorem~3.3 proves the concrete system correct.
\item If the abstract checker returns a finite or lasso-shaped counterexample,
determine symbolically, against the concrete relation \(R\), whether it has a
concrete realization.  If it has a realization, return the concrete
counterexample.
\item If it does not, properly split the failure class so that no new class
mixes a state reached by the validated concrete prefix with a state enabling
the failed continuation. Replace \(h\) by the refined abstraction and
model-check again.
\end{enumerate}
The concrete validation and separating refinements are defined below
\cite[\S4.1]{ClarkeEtAl2003}.
For a combined approximation, use the exact-recheck policy proved below
before taking a partition-refinement step.

\paragraph{Exact validation of a finite abstract path.}
Let \(n\geq1\) and let
\(\widehat T=\langle\widehat s_1,\ldots,\widehat s_n\rangle\)
be a finite abstract counterexample for a specification whose violation is
witnessed by such a path. Appropriateness of \(h\) ensures that a concrete
realization has the same specification-atom labels.
Its concrete path set is
\[
\begin{split}
 h^{-1}_{path}(\widehat T)=
 \{\langle s_1,\ldots,s_n\rangle:\;&s_1\in I,\ 
 h(s_i)=\widehat s_i\ (1\le i\le n),\\
 &R(s_i,s_{i+1})\ (1\le i<n)\}.
\end{split}
\]
For \(X\subseteq S\), define the one-step concrete forward image
\[
 \operatorname{Img}(X,R)
   =\{t\in S:\exists s\in X\;R(s,t)\}.
\]
This is a state-set operation on the concrete relation, distinct from the
inverse abstraction image \(h^{-1}\) and from the abstract relation
\(\widehat R_h\); it is needed at every validation step.  Define
\[
 S_1=h^{-1}(\widehat s_1)\cap I,
 \qquad
 S_i=\operatorname{Img}(S_{i-1},R)
       \cap h^{-1}(\widehat s_i)
 \quad(1<i\le n).
\]
Lemma~4.10 proves the equivalence of:
\[
\begin{array}{ll}
\text{(i)}&\widehat T\text{ corresponds to a concrete counterexample},\\
\text{(ii)}&h^{-1}_{path}(\widehat T)\ne\varnothing,\\
\text{(iii)}&S_i\ne\varnothing\text{ for every }1\le i\le n.
\end{array}
\]
The first two clauses are equivalent by the definition of concretization.  A
concrete path contributes its \(i\)-th state to every \(S_i\), proving
(ii)\(\Rightarrow\)(iii).  Conversely, choose \(s_n\in S_n\).  The definition
of image lets one choose \(s_{i-1}\in S_{i-1}\) with
\(R(s_{i-1},s_i)\) recursively backward to \(S_1\); the resulting sequence is
in \(h^{-1}_{path}(\widehat T)\).  This proves the criterion, not merely one
 direction of it \cite[lem.~4.10; pp.~772--773]{ClarkeEtAl2003}.

For a counterexample in the exact abstract structure,
\(\widehat s_1\in\widehat I_h\), hence
\(S_1=h^{-1}(\widehat s_1)\cap I\ne\varnothing\).  This also eliminates the
algorithm name \textsc{SplitPath}.  If
\[
 r=\min\{i\in\{2,\ldots,n\}:S_i=\varnothing\},
\]
then the procedure returns failure position \(r-1\) and
\(S_{\mathrm{prev}}=S_{r-1}\); if the set defining \(r\) is empty, the path is
concrete.  Thus \textsc{SplitPath} is precisely the recurrence plus this first-
empty-index output contract, not an additional semantic primitive
\cite[fig.~7; p.~773]{ClarkeEtAl2003}.
The combined structures defined above also use exactly \(\widehat I_h\), so
their validation starts with the same nonempty \(S_1\).

The paper's own Example~4.11 takes
\(D=\{1,\ldots,12\}\) and
\[
 h(x)=\left\lfloor\frac{x-1}{3}\right\rfloor+1.
\]
The four abstract blocks are
\[
 \{1,2,3\},\quad\{4,5,6\},\quad
 \{7,8,9\},\quad\{10,11,12\}.
\]
For the abstract path
\(\widehat T=\langle\widehat{1},\widehat{2},\widehat{3},\widehat{4}\rangle\),
Figure~6 gives
\(I\cap\{1,2,3\}=\{1,2,3\}\) and the displayed concrete arrows
\[
 1R4,\quad 1R5,\quad 3R6,\quad
 4R9,\quad 5R9,\quad 7R12.
\]
The recurrence is therefore reproducible:
\[
\begin{aligned}
 S_1&=I\cap\{1,2,3\}=\{1,2,3\},\\
 S_2&=\operatorname{Img}(S_1,R)\cap\{4,5,6\}
       =\{4,5,6\},\\
 S_3&=\operatorname{Img}(S_2,R)\cap\{7,8,9\}
       =\{9\},\\
 S_4&=\operatorname{Img}(\{9\},R)\cap\{10,11,12\}
       =\varnothing.
\end{aligned}
\]
Thus \(\widehat{3}\), the abstract state at the last index whose propagated
reachable set \(S_i\) is nonempty, is the failure state.  The otherwise
unreachable state \(7\) supplies the misleading abstract edge through
\(7R12\), while the reached state \(9\) cannot continue.  The source's
symbolic procedure computes this first empty step and the preceding set; it
does not infer spuriousness from the abstract edges alone
\cite[ex.~4.11 and fig.~6]{ClarkeEtAl2003}.

\paragraph{What the refinement must separate.}
At a finite-path failure between \(\widehat s_i\) and
\(\widehat s_{i+1}\), the source partitions the concrete states represented by
\(\widehat s_i\) into
\[
\begin{aligned}
 S_D&=S_i,\\
 S_B&=\{s\in h^{-1}(\widehat s_i):
       \exists t\in h^{-1}(\widehat s_{i+1})\ R(s,t)\},\\
 S_I&=h^{-1}(\widehat s_i)\setminus(S_D\cup S_B).
\end{aligned}
\]
The standing meanings and distinctions are now exact.  States in \(S_D\) are
reachable along the validated prefix but have no continuation into the next
abstract block.  States in \(S_B\) supply the misleading abstract transition
but were not reached by that prefix. States in \(S_I\) do neither.
The separating refinement used here requires that no new abstract class
contain both an \(S_D\)-state and an \(S_B\)-state. This disconnects the
validated concrete prefix from the failed continuation within the new class.
It guarantees a proper partition refinement; a subsequent abstract check
may still find a different infeasible path with the same old coarse
projection.

In Example~4.11, these three parts of the failure block are
\[
 S_D=\{9\},\qquad S_B=\{7\},\qquad S_I=\{8\}.
\]
Splitting \(9\) from \(7\) removes the counterexample: the refined block reached
by the prefix has no edge to \(\widehat4\), whereas the refined block retaining
that edge is not reached by the prefix.  The state \(8\) may be placed on
either side without changing this argument.  Model checking then resumes on
the refined abstraction \cite[ex.~4.15 and \S4.4]{ClarkeEtAl2003}.

For the composite \(h\) defined above, a componentwise refinement consists of
equivalences \(\equiv'_j\) satisfying
\[
 =_{D_j}\ \subseteq\ \equiv'_j\ \subseteq\ \equiv_j
 \qquad(1\le j\le m),
\]
with \(\equiv'_h\) again their product. It is separating when no
\(d\in S_D\) and \(b\in S_B\) remain \(\equiv'_h\)-equivalent. Write the
failure class as
\[
 C=h^{-1}(\widehat s_i)=E_1\times\cdots\times E_m,
\]
where each \(E_j\) is one old \(\equiv_j\)-class. The size of a refinement
at this failure is
\[
 \bigl|C/{\equiv'_h}\bigr|
   =\prod_{j=1}^{m}\bigl|E_j/{\equiv'_j}\bigr|.
\]
Thus minimizing size means minimizing the number of rectangular product
blocks replacing \(C\). A coarsest refinement in the order of relation
inclusion would be stronger: it would contain every separating refinement.
Such a greatest relation need not exist when \(S_I\ne\varnothing\).
Finding a minimum-size separating product refinement is NP-hard
\cite[thm.~4.17 and \S5]{ClarkeEtAl2003}.

\paragraph{Why minimum-size product refinement is hard.}
A finite matrix with entries \(0,1,*\) represents a two-coordinate failure
class: rows and columns are the coordinate domains, \(0\)-cells form \(S_D\),
\(1\)-cells form \(S_B\), and \(*\)-cells form \(S_I\). Initially each
coordinate has one equivalence class. A pair of row and column partitions
is separating precisely when none of its rectangular blocks contains both
a \(0\) and a \(1\); the symbol \(*\) imposes no constraint. Its size is
the number of row classes times the number of column classes.

First reduce clique partition to column partition with all rows fixed.
For an undirected graph \(G=(\{1,\ldots,n\},E)\), \(n\geq1\), and a
threshold \(1\leq K\leq n\), form
\[
 M_{ij}=
 \begin{cases}
  1,&i=j,\\
  0,&i\ne j\text{ and }\{i,j\}\notin E,\\
  *,&i\ne j\text{ and }\{i,j\}\in E.
 \end{cases}
\]
Two columns are compatible if no row places a \(0\) in one and a \(1\)
in the other. Distinct columns \(i,j\) are compatible exactly when
\(\{i,j\}\in E\). Indeed, their only possible conflicts occur in rows
\(i,j\), since those are the only rows containing a \(1\) in either
column. A group of columns is compatible exactly when all its pairs are
compatible: any forbidden \(0/1\) pair already lies in two columns.
Consequently a column group is compatible exactly when its vertices form
a clique. The graph admits a partition into at most \(K\) cliques if and
only if \(M\) admits at most \(K\) compatible column groups. Clique
partition is the graph-colouring problem on the complement graph, so this
is an NP-hard decision problem.

We must still prevent row partitions from making the size problem easier.
The following construction does so for any \(n\times p\) column instance
\(M\), with \(n,p\geq1\). Put
\[
 t=\lceil\log_2 n\rceil,\qquad q=t+1.
\]
Append \(q\) columns whose entries in original row \(i\) are
\[
 \bigl(1,\operatorname{bin}_t(i-1)\bigr),
\]
where \(\operatorname{bin}_t(a)\) is the \(t\)-bit binary representation
of \(a\), including leading zeros. Append one further row with \(1\) in
every original column and \(0\) in every appended column. Denote the
resulting \((n+1)\times(p+q)\) matrix by \(B(M)\).
Its additional row and columns have four properties.
\begin{enumerate}
\item Distinct original rows conflict in an appended binary column, because
their identifiers differ.
\item Every original row conflicts with the additional row in the first
appended column.
\item Every appended column conflicts with every original column in the
additional row.
\item Appended columns conflict pairwise. The leading column differs from
each binary column in row \(1\), whose identifier is zero. Two distinct
binary columns differ in a row with identifier \(2^r\) for one of their
bit positions \(r\). This row exists: for \(n>1\),
\(2^r\leq2^{t-1}\leq n-1\). For \(n=1\), there are no binary columns,
and the single leading column suffices.
\end{enumerate}
A conflicting pair cannot belong to a common row or column class in any
separating product partition, irrespective of what happens in the other
coordinate: its two conflicting cells would then lie in the same
rectangle. Thus all \(n+1\) rows remain separate, all \(q\) appended
columns remain separate, and an appended column cannot join an original
column. The additional row gives equal entries to the original columns,
so it introduces no new compatibility constraint between them. The only
remaining choice is the compatible partition of the original columns.
It follows that
\[
 \begin{split}
 M\text{ has at most }K\text{ compatible column groups}
 \quad\Longleftrightarrow\quad
 B(M)\text{ has a separating product partition}\\
 \text{of size at most }(n+1)(K+q).
 \end{split}
\]
Both the construction and its threshold have polynomial size. Interpreting
the entries of \(B(M)\) as \(S_D,S_B,S_I\) completes the reduction to
minimum-size separating product refinement, already for two coordinate
domains given by explicit finite tables
\cite[thm.~4.17 and \S5]{ClarkeEtAl2003}.

\paragraph{A coarsest refinement for separation from the complement.}
There is a simpler problem with an explicit solution: separate
\(A\subseteq C\) from all of \(C\setminus A\), rather than leaving the
placement of \(S_I\) unconstrained. Assume \(m\geq1\).

The polynomial operation needs one further typed map.  For
\(X\subseteq D_1\times\cdots\times D_m\), \(a\in D_j\), define
\[
\begin{split}
 \operatorname{proj}(X,j,a)=
 \{(d_1,\ldots,d_{j-1},d_{j+1},\ldots,d_m):\;&\\
 (d_1,\ldots,d_{j-1},a,d_{j+1},\ldots,d_m)&\in X\}.
\end{split}
\]
This projection deletes coordinate \(j\) after fixing its value; it is
different from the variable-cluster coordinate restriction
\(\pi_{VC_j}\) above and is used only to compare slices of a failure block.
For \(m=1\), the product over the remaining coordinates is the singleton
\(\{()\}\), so a slice is either empty or \(\{()\}\).
Define the component update
\[
 a\equiv'_j b
 \iff
 a\equiv_j b
 \ \land\
 \operatorname{proj}(A,j,a)
   =\operatorname{proj}(A,j,b)
 \qquad(1\le j\le m).
\]
Its product is the unique greatest product refinement separating \(A\)
from \(C\setminus A\).

\emph{Proof.}
Equality of slices is an equivalence relation; intersecting it with
\(\equiv_j\) gives an equivalence contained in \(\equiv_j\). Hence the
product update is an equivalence refining the old one.
For separation, suppose \(x\in A\) and \(x\equiv'_h y\). Replace the
coordinates of \(x\) successively:
\[
 x^{(0)}=x,\qquad
 x^{(j)}=(y_1,\ldots,y_j,x_{j+1},\ldots,x_m).
\]
At step \(j\), equality of the slices for \(x_j,y_j\), evaluated at
the other currently fixed coordinates, gives
\[
 x^{(j-1)}\in A\quad\Longleftrightarrow\quad x^{(j)}\in A.
\]
Induction gives \(y\in A\), proving separation.

Now let \(\approx_h\) be any other separating product refinement with
component relations \(\approx_j\subseteq\equiv_j\). Suppose
\(a\approx_j b\) but their \(A\)-slices differ. After interchanging
\(a,b\) if necessary, take
\[
 z\in\operatorname{proj}(A,j,a)
       \setminus\operatorname{proj}(A,j,b).
\]
Insert \(a\), respectively \(b\), into coordinate \(j\) of \(z\) to form
full states \(x,y\). Then \(x\in A\subseteq C\). Since \(C\) is the
rectangle \(\prod_\ell E_\ell\), all other coordinates of \(z\) lie in
their \(E_\ell\), and \(a\in E_j\). The implication
\(a\approx_j b\Rightarrow a\equiv_j b\) puts \(b\) in the same \(E_j\);
thus \(y\in C\setminus A\). But \(x\approx_h y\), using
\(a\approx_j b\) in coordinate \(j\) and reflexivity elsewhere, contrary
to separation. Therefore every pair in \(\approx_j\) has equal slices:
\[
 \approx_j\subseteq\equiv'_j\quad\text{for every }j,
 \qquad
 \approx_h\subseteq\equiv'_h.
\]
The update is itself separating and contains every separating product
refinement, which proves greatestness and uniqueness. \(\square\)

Because \(C\) is finite, every proper refinement of its partition has
strictly more blocks. The greatestness just proved therefore also gives
the unique minimum partition on \(C\). The coordinate update can also
split other old product classes that use the same coordinate values.
Additional splits beyond this update cannot improve the local size
objective; uniqueness of the minimum on \(C\) does not by itself
prohibit extra global splits outside \(C\).

Apply the update with \(A=S_D\). This is the operation
\textsc{PolyRefine}. When \(S_I=\varnothing\), its complement in \(C\)
is exactly \(S_B\), so the result is the coarsest separating refinement
and has minimum size. When \(S_I\ne\varnothing\), it solves the
stronger problem of separating \(S_D\) from \(S_B\cup S_I\).
It still separates \(S_D,S_B\), but need not minimize the size required
by that weaker condition
\cite[lems.~4.18--4.19, cors.~4.20, 4.22, and rem.~4.23;
pp.~779--781]{ClarkeEtAl2003}.

For example, start with universal coordinate equivalences on
\(C=\{0,1\}^2\) and take
\[
 S_D=\{(0,0)\},\qquad S_B=\{(1,1)\},\qquad
 S_I=\{(0,1),(1,0)\}.
\]
Splitting only the first coordinate separates \(S_D,S_B\) with two
blocks; splitting only the second does the same. These two equivalences
are incomparable. Any product equivalence containing both must be
universal in each coordinate and hence fails separation, so the weaker
problem has no greatest separating refinement. The slice update splits
both coordinates, giving four blocks. This is necessary for separating
\(S_D\) from its entire complement, but not for separating it from
\(S_B\) alone.

The update is polynomial when the coordinate domains and state sets are
given as explicit finite tables: form each slice, compare these finite
tables, and group equal slices within each old coordinate class. This
statement is about the explicit representation size, not the number of
Boolean variables in a succinct encoding and not the size of arbitrary
intermediate BDDs.

At an exact abstract failure, \(S_D\ne\varnothing\) by choice of the
first empty propagation step. Also \(S_B\ne\varnothing\), since the
failed abstract edge has a concrete witness in the exact existential
image. The two sets are disjoint, or the propagation would not have
failed. Any separating update therefore separates two states formerly
in the same \(C\), so it properly refines the global state partition.

\paragraph{Validation of an abstract lasso.}
For \(n\geq1\) and \(0\leq i<n\), write an abstract lasso as
\[
 \widehat T=
 \langle\widehat s_1,\ldots,\widehat s_i\rangle
 \langle\widehat s_{i+1},\ldots,\widehat s_n\rangle^\omega
\]
and put
\[
 m=\min_{i+1\le j\le n}
      \lvert h^{-1}(\widehat s_j)\rvert.
\]
Let \(\widehat T^{unwind}\) be the finite abstract path consisting of the
whole prefix followed by \(m+1\) copies of the loop segment:
\[
 \widehat T^{unwind}
 =
 \langle\widehat s_1,\ldots,\widehat s_i\rangle
 \langle\widehat s_{i+1},\ldots,\widehat s_n\rangle^{m+1}.
\]
Theorem~4.13 proves
\[
 \widehat T\text{ has a concrete lasso realization}
 \iff
 h^{-1}_{path}(\widehat T^{unwind})\ne\varnothing.
\]
The forward implication takes a finite prefix of a concrete infinite
realization.  For the reverse implication, choose
\[
 M\in\{i+1,\ldots,n\}
 \quad\text{with}\quad
 \lvert h^{-1}(\widehat s_M)\rvert=m.
\]
The unwound concrete path visits that same loop coordinate once in each of the
\(m+1\) whole loop copies, so two occurrences contain the same concrete
state.  Their positions differ by an integral number of complete loop
traversals.  The concrete segment between them is therefore a cycle projecting
to complete copies of the abstract loop; repeating it produces a concrete
lasso of the required abstract shape.  This pigeonhole argument is why a
bounded number of unwindings suffices even though the eventual period of the
reachable-set sequence can be exponential.

The source's \textsc{SplitLoop} name can also be eliminated.  For a position
\(q\) in the unwound word, define its abstract loop index by
\[
 \lambda(q)=
 \begin{cases}
  q,&q\le n,\\
  i+1+\bigl((q-i-1)\bmod(n-i)\bigr),&q>n.
 \end{cases}
\]
Run the same concrete-image recurrence along
\(\widehat T^{unwind}\).  If \(r\) is its first empty index, the procedure
returns
\[
 S_{\mathrm{prev}}=S_{r-1},\qquad
 p=\lambda(r-1),\qquad k=\lambda(r),
\]
and these outputs satisfy
\[
\begin{aligned}
 S_{\mathrm{prev}}&\subseteq h^{-1}(\widehat s_p),\\
 S_{\mathrm{prev}}&\text{ is reachable from }
   h^{-1}(\widehat s_1)\cap I\text{ along the unwound prefix},\\
 k&=p+1\quad(p<n),\qquad k=i+1\quad(p=n),\\
 \operatorname{Img}(S_{\mathrm{prev}},R)
   \cap h^{-1}(\widehat s_k)&=\varnothing.
\end{aligned}
\]
Hence \(\widehat s_p\) is the failure state.  Refinement takes
\(S_D=S_{\mathrm{prev}}\) and computes \(S_B,S_I\) against the successor block
\(\widehat s_k\); the name does not conceal an additional operation
\cite[ex.~4.12, thm.~4.13, and fig.~9; pp.~774--776]{ClarkeEtAl2003}.

\paragraph{Termination and completeness of exact refinement.}
Let \(M\) be finite, with nonempty state set \(S\), and let \(\varphi\)
be an ACTL* specification in a fragment for which a failed model check
returns a finite-path or lasso counterexample. The counterexample
representation must witness violation: every concrete realization with
the same specification-atom labels is a concrete refutation.
Assume the model checker and the finite symbolic set operations terminate.
Then the exact refinement loop terminates with a correct verdict and an
abstract model \(\widehat M_h\) satisfying
\[
 \widehat M_h\models\varphi
 \quad\Longleftrightarrow\quad
 M\models\varphi.
\]

\emph{Proof.}
The initial product abstraction is appropriate. A refinement never
identifies states that were previously distinct, so it cannot identify
states disagreeing on a specification atom. Appropriateness is therefore
invariant. The simulation-preservation theorem consequently applies at
every iteration.

Put \(N=|S|\) and let \(b(h)=|S/{\equiv_h}|\) be the number of current
abstract state classes. At a spurious finite counterexample, validation
finds a first empty set after a nonempty one. At a spurious lasso, the
finite unwinding criterion gives exactly the same kind of failure.
The exact initial image rules out an empty first set. Since the failed
edge belongs to the exact abstract relation, both \(S_D\) and \(S_B\)
are nonempty, disjoint subsets of one old class. The separating product
update just proved is therefore proper: it increases \(b(h)\) and
never merges any other classes. Thus \(N-b(h)\) is a nonnegative
integer that strictly decreases at every refinement.

There can be at most \(N-b(h_0)\) refinements, where \(h_0\) is the
initial abstraction. No spurious counterexample can require a further
split once \(b(h)=N\): every class is then a singleton, and the exact
abstract structure is isomorphic to \(M\). Equivalently, the two
nonempty disjoint sets needed for a failure could not both lie in a
singleton. Hence the loop must reach a terminating branch.

If an abstract check succeeds, simulation proves \(M\models\varphi\).
If validation succeeds, its finite-path or lasso construction supplies
a concrete counterexample, proving \(M\not\models\varphi\). In that
branch the current abstract check also failed. Both terminating branches
therefore give the stated equivalence. \(\square\)

This establishes the finite-state completeness assertion for the
path-or-lasso counterexample fragment
\cite[thm.~4.24]{ClarkeEtAl2003}. A specification whose refutation
requires a branching counterexample is outside these hypotheses; its
validator and progress argument would have to treat that branching
object.

\paragraph{A terminating policy for combined approximation.}
The combined relation can contain an edge with no concrete witness at all.
At such an edge the validation set \(S_B\) may be empty, so a requested
separation from \(S_D\) would not force a proper partition refinement.
Use the following exact-recheck policy with
\(\widetilde M_{h,J}\), retaining its exact initial image.
\begin{enumerate}
\item Check \(\widetilde M_{h,J}\). If it satisfies \(\varphi\),
return success. Otherwise validate its finite-path or lasso
counterexample against \(R\). If validation succeeds, return the
concrete counterexample.
\item If validation fails and \(J\ne\{1,\ldots,q\}\), set
\(J=\{1,\ldots,q\}\) and check again at the same \(h\).
This check uses \(\widehat R_h\). Keep this exact selection until
the current partition changes; the exact recheck may succeed or may
produce a different counterexample.
\item If validation fails with \(J=\{1,\ldots,q\}\), apply the proper
separating product refinement. At the new \(h\), a new selection \(J\)
may be chosen, and the process repeats.
\end{enumerate}
In particular, a partition split is attempted only after validation of
a counterexample in the exact abstract structure. The method can use
early abstraction at each partition without depending on it to supply
the witnesses needed for partition progress
\cite[\S6]{ClarkeEtAl2003}.

\emph{Termination and correctness.}
Under the same finite-state and counterexample hypotheses, define
\[
 \epsilon(J)=
 \begin{cases}
  1,&J\ne\{1,\ldots,q\},\\
  0,&J=\{1,\ldots,q\}.
 \end{cases}
\]
Order pairs of nonnegative integers lexicographically, comparing the
first coordinate first. At every nonterminating update the pair
\[
 \bigl(N-b(h),\epsilon(J)\bigr)
\]
strictly decreases. An exact recheck keeps \(h\) fixed and changes
\(\epsilon\) from \(1\) to \(0\). An exact spurious counterexample
properly increases \(b(h)\), decreasing the first coordinate even if a
new approximate selection resets \(\epsilon\) to \(1\). Since the
second coordinate is always \(0\) or \(1\), there is no infinite
descending sequence, so the policy terminates.

At every selection,
\(M\preceq\widehat M_h\preceq\widetilde M_{h,J}\); a successful
combined check therefore proves the concrete specification. A validated
counterexample proves its concrete failure by the same realization
criterion as before. At the identity partition an exact recheck has
neither extra transitions nor extra initial states, so no nonterminating
case remains. These facts prove completeness of this specified combined
policy. They require no assumption that an arbitrary unsuccessful
approximation heuristic will eventually improve itself. \(\square\)

\paragraph{The industrial experiment.}
On a 200\,MHz
PentiumPro with 1\,GB RAM, the Fujitsu case began from about 61,500 lines of
Verilog; after engineers removed irrelevant units and abstracted the datapath,
the verification model had about 500 latches and was translated into about
9,500 lines of SMV.  The method abstracted 144 symbolic variables, used three
refinement steps, and exposed a previously undiscovered bug after the compared
NuSMV and Yang SMV runs failed to verify the design.  The authors had added
cone-of-influence reduction to NuSMV but not to Yang's SMV, so those two
comparison configurations were not symmetric.
Those numbers report one industrial experiment; they are not a general
complexity guarantee \cite[\S7.2]{ClarkeEtAl2003}.
